\documentclass[11pt]{article}
\usepackage[T1]{fontenc}\usepackage[utf8]{inputenc}
\usepackage{lmodern,amsmath,amssymb,amsthm,mathtools,microtype,booktabs,array}
\usepackage[margin=1in]{geometry}
\usepackage[authoryear,round]{natbib}
\usepackage[colorlinks=true,allcolors=blue]{hyperref}
\newcommand{\workbenchurl}{https://github.com/KokunoYumeto/schemes-and-moduli-workbench}
\newcommand{\recordurl}{https://github.com/KokunoYumeto/mathematics-workbenches}
\newcommand{\Z}{\mathbb Z}\newcommand{\N}{\mathbb N}\newcommand{\Q}{\mathbb Q}
\newcommand{\A}{\mathbb A}\newcommand{\B}{\mathbb B}\newcommand{\F}{\mathbb F}\newcommand{\C}{\mathbb C}
\newcommand{\eps}{\varepsilon}\newcommand{\supp}{\operatorname{supp}}
\newcommand{\Spec}{\operatorname{Spec}}\newcommand{\tr}{\operatorname{tr}}
\newtheorem{theorem}{Theorem}[section]\newtheorem{proposition}[theorem]{Proposition}
\newtheorem{lemma}[theorem]{Lemma}\newtheorem{corollary}[theorem]{Corollary}
\theoremstyle{definition}\newtheorem{definition}[theorem]{Definition}
\newtheorem{remark}[theorem]{Remark}\newtheorem{example}[theorem]{Example}
\title{Effective and ineffective lifts of Mulmuley's\\rank-three algebra, after Forsberg}
\author{}\date{24 November 2026}
\begin{document}\maketitle\enlargethispage{\baselineskip}
\begin{abstract}
Faul, Goswami, Joubert, and Manuell's zero adjunction separates an absent coordinate from a present coordinate of value zero, while Mulmuley's nonstandard canonical-basis programme asks for positive and unimodal structure constants. We construct an associative inverse-semiring lift of the rank-three table with no supported zero in any bottom coefficient, and prove that preserving the bottom ideal makes this impossible. We also describe the restricted adelic semiring and its units. Five effective associative lifts are reconstructed by matrix multiplication; two have unimodal dimension polynomials. One lift has only 38 occupied zero-trace weights. The zero-trace weights are counted separately from cancelling multiplicities. A displayed coefficient in Mulmuley's rank-three table has a nonunimodal positive factor, contrary to the verification claimed there. Applying Forsberg's obstruction, we construct a second associative support lift that admits no effective dimension lift. At every rank we also specify a finite torsion-free integral word order over $\mathbb Z[D]$, with faithful Laurent base change and a localization certificate. Global minimization and the least positive-basis coefficient ring at rank four remain open. The mathematics and text were produced by GPT-6 Astra Pro and separately checked by the same model; no human has checked them.
\end{abstract}

\section{The arithmetic support and the lifting question}
For a commutative ring $R$, put $R_0=R\sqcup\{\tau\}$ with the new additive zero and multiplicative absorber $\tau$. The old ring zero is denoted $e$. This construction is Definition 3.6 of \citet{FGJM}; their Proposition 6.1 gives the inverse additive structure of its semimodules. The data of an idempotent semiring and an abelian-group-valued lax monoidal functor describe inverse semirings by \citet[Theorem 8.5]{CFMM}. We use these descriptions to separate support, signed trace, and effective multiplicity.

The first theorem records exactly which familiar adelic object survives zero adjunction. A place means either a finite prime or the real place of $\Q$.

\begin{theorem}[Restricted adelic support]\label{thm:adeles-main}
The restricted product
\[
 \mathcal A_0=\prod_v'(\Q_v)_0
 \quad\hbox{with respect to }(\Z_p)_0\hbox{ at finite places}
\]
is the semiring of pairs $(S,a)$, where $S$ is any subset of the places, $a\in\A_\Q$, and $a_v=0$ outside $S$. The operations are
\[
 (S,a)+(T,b)=(S\cup T,a+b),\qquad
 (S,a)(T,b)=(S\cap T,ab).
\]
Its units are exactly the ordinary ideles, with support all places. If
\[
 \sum_{v\in S}\log|r|_v=0\quad\hbox{for every }r\in\Q^\times,
\]
then $S$ is empty or is the set of all places.
\end{theorem}

The condition in the last assertion is a product-formula condition on a support. It makes precise why selecting only some ordinary places is not a new balanced arithmetic curve. It does not exclude additional trivial absolute values, whose logarithms are identically zero. The unit statement also corrects the claim that $(\Q_v)_0$ is a semifield: its old zero is nonzero as a semiring element and is not invertible.

For the lifting theorems let $q$ be an indeterminate, $D=q+2+q^{-1}$, and write
\[
 \mathcal E=\N[q,q^{-1},\eps]/(\eps^2-1),\qquad
 \mathcal S=\Z_0[q,q^{-1}].
\]
The map $\pi:\mathcal E\to\mathcal S$ records the support of the dimension and the signed trace. A zero-trace weight is \emph{supported} when it has positive dimension. The bottom basis vector of the rank-three algebra will be denoted $\Sigma$.

\begin{theorem}[An unrestricted lift with no bottom supported zeros]\label{thm:no-bottom-main}
Let $B=\{1,\Sigma\}\cup\{\gamma_w^r:r=1,2,\ w\in\{1,2,12,21\}\}$ be the basis of the rank-three trace algebra defined in Section~\ref{sec:rank3}. There is an associative unital algebra on the free $\mathcal S$-semimodule with basis $B$ whose coefficient traces are precisely that table and whose bottom coefficients have exactly their nonzero trace supports. In particular no bottom coefficient has a supported zero-trace weight. Its span of $\Sigma$ is not an ideal. An explicit construction is given in Section~\ref{sec:unrestricted}.
\end{theorem}

Thus the coefficientwise trace-lifting condition alone does not force a supported zero into the bottom cell. This resolves the unrestricted inverse-semiring version of a proposed obstruction in the project's unpublished lifting notes. It is distinct from an assertion about cancelling multiplicities in effective dimension rings: a cancelling pair at a nonzero trace weight is not itself a supported zero. The following theorem identifies the precise ideal-preserving obstruction.

\begin{theorem}[Bottom-ideal obstruction]\label{thm:bottom-main}
Let a finite free $\mathcal S$-semimodule with a named basis carry an associative unital algebra structure. Suppose the basis contains $\gamma$ and $\Sigma$, its trace specialization has
\[
 c(\gamma,\gamma;\gamma)=-D(D-1)(D-4),\qquad
 \gamma\Sigma=\Sigma\gamma=4D(D-1)\Sigma,
\]
and the lifted span of $\Sigma$ is a two-sided ideal. If $\gamma\Sigma=g\Sigma$ in the lift, then
\[
 \max\supp g\ge3,\qquad \min\supp g\le-3.
\]
Since the trace of $g$ has support exactly $\{-2,-1,0,1,2\}$, at least two supported zero-trace weights occur in $g$.
\end{theorem}

This theorem is stronger than an assertion that some positive and negative multiplicities cancel. Its proof uses the stated bottom-ideal hypothesis. A general lift may add coefficients with zero trace in other outputs of $\gamma\Sigma$. Theorem~\ref{thm:no-bottom-main} shows that removing the ideal hypothesis allows the conclusion to fail.

\begin{theorem}[Effective rank-three lifts]\label{thm:lifts-main}
The ten-dimensional rank-three algebra described in Section~\ref{sec:rank3} has associative unital lifts of its structure constants to $\mathcal E$ preserving the pattern of nonzero basis outputs. Among five explicit lifts, one has $38$ supported zero-trace weights and $8{,}288$ cancelling pairs. Another has $42$ supported zero-trace weights and $10{,}804$ cancelling pairs, with all $132$ nonzero dimension polynomials symmetric and unimodal. Each has $34$ supported zero-trace weights in bottom coefficients. A third has $7{,}656$ cancelling pairs and $54$ supported zero-trace weights. These are counts for specified lifts; no global minimum is asserted.
\end{theorem}

Mulmuley's positivity conjecture asks for more than the existence of such lifts. In particular it asks for unimodality of a positive factor of the trace coefficient itself \citep[Conjecture 3.4, p.\ 24]{Mulmuley}. A coefficient printed in his Figure 2 has positive factor
\[
 (q+1+q^{-1})(q+q^{-1})^2
 =q^3+q^2+3q+2+3q^{-1}+q^{-2}+q^{-3}.
\]
Its middle dip disproves the corresponding assertion in the rank-three verification on p.\ 29, under that displayed normalization. Replacing trace coefficients by dimension coefficients is a different statement; the two unimodal lifts prove that statement for the multiplication table defined here.


\begin{theorem}[Failure of effective lifting, after Forsberg]\label{thm:nonliftable-main}
On the same rank-three basis there is an associative unital $\mathcal S$-algebra with the prescribed trace table that is not the image under $\pi$ of any associative unital $\mathcal E$-algebra on that basis. More strongly, its occupied basis-output relation admits no associative choice of finite positive real multiplicities.
\end{theorem}

Forsberg's multisemigroup obstruction shows that an idempotent output relation need not lift to finite multiplicities \citep[Proposition 12 and Corollary 13]{Forsberg}. Here we place that obstruction inside an associative inverse-semiring lift of the specified rank-three algebra. Thus the coefficientwise map in the dimension criterion is surjective, but its induced map on associative structures is not. The construction and the full specialized proof are in Section~\ref{sec:nonliftable}.

\begin{theorem}[Least free basis rank for an obstruction]\label{thm:minimal-basis-rank-main}
Among associative unital algebra structures on a finite free $\mathcal S$-semimodule with a named basis containing the strict unit, every structure of basis rank one or two is the coefficientwise image of an associative unital $\mathcal E$-algebra on the same basis. There is a commutative associative structure of basis rank three with no such refinement, whose ordinary trace algebra is
\[
 \mathbb Z[q,q^{-1}][a,b]/(a,b)^2.
\]
Thus three is the least free basis rank at which coefficientwise surjectivity fails to lift an associative algebra structure.
\end{theorem}

Here basis rank counts the named vectors, including the unit. It is distinct from the nonstandard algebra's rank parameter: its rank-three table has ten named vectors. The obstruction uses the same two-nonunit output relation as Forsberg's example; the lower-rank lifting argument fixes the strict-unit convention explicitly.

The next theorem specifies an integral coefficient order for the generators of \citet[Definition 11.1, equation (93), p.~59]{BSM2013}. It does not assume a positive or canonical basis.

Let $v$ be an indeterminate and put
\[
 D=(v+v^{-1})^2,\qquad S_D=\mathbb Z[D],\qquad
 R_v=\mathbb Z[v,v^{-1}],\qquad
 F_D=\mathbb Q(D),\qquad K_v=\mathbb Q(v).
\]
Thus the earlier Laurent parameter is $q=v^2$. For a positive integer $r$, let $H_{R_v}(\mathfrak S_r)$ be the type-A Hecke algebra with its standard basis $T_w$ and relations
\[
 (T_s-v)(T_s+v^{-1})=0
\]
for the simple reflections, together with the braid relations. In
\[
 M_{R_v}=H_{R_v}(\mathfrak S_r)\otimes_{R_v}
 H_{R_v}(\mathfrak S_r)
\]
define
\[
 P_s=(T_s+v^{-1})\otimes(T_s+v^{-1})
       +(T_s-v)\otimes(T_s-v).
\]
A \emph{word} means a finite product of the $P_s$, including the empty word $1$. Write $A_{S_D},A_{R_v},A_{F_D},A_{K_v}$ for the spans of these words over the indicated rings or fields, all viewed in the ambient Hecke tensor algebra over $K_v$.

\begin{theorem}\label{thm:hecke-word-descent}
For every positive integer $r$, the following statements hold.
\begin{enumerate}
\item The algebra $A_{S_D}$ is a finite torsion-free $S_D$-module.
\item Multiplication induces isomorphisms of algebras
\[
 R_v\otimes_{S_D}A_{S_D}\xrightarrow{\ \sim\ }A_{R_v},
 \qquad
 K_v\otimes_{F_D}A_{F_D}\xrightarrow{\ \sim\ }A_{K_v}.
\]
\item The algebra $A_{F_D}$ has an $F_D$-basis consisting of words. Every multiplication coefficient in such a basis belongs to $F_D$.
\item For $a\in A_{S_D}$, let $L_a$ be left multiplication on the finite-dimensional $F_D$-vector space $A_{F_D}$. Then
\[
 \det(X-L_a)\in S_D[X].
\]
In particular, the regular trace of every word lies in $\mathbb Z[D]$.
\end{enumerate}
\end{theorem}



\section{Proof overview}
The restricted adelic product can be read coordinate by coordinate. A new zero contributes no support, while an old zero contributes a supported zero in the ordinary adele. Multiplication intersects supports, and a unit therefore has every place present. Testing the product formula on one rational prime at a time gives the support rigidity over $\Q$.

The lifting argument uses two independent evaluations of the sign, at $+1$ and $-1$. These identify the coefficient ring with pairs of Laurent polynomials congruent modulo two. Associativity of an effective lift is therefore exactly associativity of its dimension and trace tables. We realize both tables in a product of two scalar algebras and two $2$ by $2$ matrix algebras. This proves associativity before checking positivity. A small finite list of Laurent coefficient inequalities then checks effectivity.

For the supported-zero obstruction, no enumeration is needed. Associativity and the bottom-ideal condition put the support of a known cubic Laurent coefficient times $g$ inside the support of $g^2$. Maximum and minimum exponents force the support of $g$ beyond the support of its trace. For the converse example, we keep each bottom support unchanged and enlarge all other outputs to centered exponent intervals. A graded radius formula makes both associations of every triple have the same interval. Their traces already agree, so the new coefficient algebra is associative. A second support construction uses one singleton output and otherwise all nonunit outputs. Its support equations are automatic, while the coefficient of a single basis vector in a triple product would acquire an extra strictly positive term in any effective lift. This is the specialization of Forsberg's obstruction.

The integral word-order argument uses the two semilinear involutions of the Hecke algebra to descend word coordinates to $\mathbb Q(D)$. Finiteness follows inside a finite free $\mathbb Z[D]$-module. Flatness then gives integral base change, and normality puts regular characteristic coefficients in $\mathbb Z[D]$. A selected nonzero coordinate minor supplies a denominator bound by Cramer's rule.

\section{Adeles, support, and the product formula}\label{sec:adeles}
This section proves Theorem~\ref{thm:adeles-main} and gives a standard geometric interpretation of its support component.

\begin{proof}[Proof of Theorem~\ref{thm:adeles-main}]
Given $x=(x_v)$ in the restricted product, set $S=\{v:x_v\ne\tau\}$ and replace every $\tau$ by the ordinary zero to obtain $a$. The restricted integrality condition is exactly the adele condition on $a$. Conversely $(S,a)$ determines $x_v=a_v$ at a place of $S$ and $x_v=\tau$ elsewhere, even when $a_v=0$ at a place of $S$. These operations are inverse. A sum is supported precisely when one summand is supported; a product is supported precisely when both are. This proves the formulas.

The multiplicative identity has every place present and value one. An invertible pair must have every place present, and its amplitude and the coordinatewise inverse amplitude must both be adeles. This is exactly the definition of an idele. Conversely an idele and its inverse give inverse elements of the restricted product.

Finally, evaluate the displayed balance identity on a positive rational prime $p$. Its only nonzero logarithms are $\log|p|_\infty=\log p$ and $\log|p|_p=-\log p$. Therefore $1_S(\infty)=1_S(p)$ for every $p$. This says that every place is in $S$, or none is.
\end{proof}

\begin{proposition}[Real weights on the rational places]\label{prop:weights}
Let $(w_v)$ be arbitrary real numbers. The expression
\[
 \ell_w(r)=\sum_v w_v\log|r|_v
\]
is a finite sum for each $r\in\Q^\times$. It vanishes for every $r$ if and only if all $w_v$ are equal. Hence a support-selected logarithmic height descends through multiplication by every rational scalar exactly for the two supports in Theorem~\ref{thm:adeles-main}.
\end{proposition}
\begin{proof}
The prime test gives $(w_\infty-w_p)\log p=0$, hence $w_p=w_\infty$. The converse is the rational product formula, obtained by writing $r=\pm\prod_pp^{v_p(r)}$. A change of homogeneous coordinates by $r$ adds exactly $\ell_w(r)$ to the weighted logarithmic expression, proving the last assertion.
\end{proof}

% Ready-to-input section; bibliography entries are supplied in the companion note.
\subsection{Real weights over global fields}
For a number field $K$, use one place for each real embedding and each
conjugate pair of complex embeddings, and set
\[
 L_v(a)=
 \begin{cases}
  \log|\sigma_v(a)|,&v\text{ real},\\
  2\log|\sigma_v(a)|,&v\text{ complex},\\
  -\operatorname{ord}_v(a)\log N(v),&v\text{ finite}.
 \end{cases}
\]
For a global function field with full constant field $\F_q$, set
$L_v(a)=-\deg(v)\operatorname{ord}_v(a)\log q$.
These conventions absorb the usual local weights into $L_v$; the product
formula is $\sum_vL_v(a)=0$.

\begin{theorem}[Rigidity of real global weights]\label{thm:global-weights}
Let $K$ be a global field and $(w_v)_v$ any family of real numbers indexed
by its standard places. Then
\[
 \sum_v w_vL_v(a)=0\quad\text{for every }a\in K^\times
 \quad\Longleftrightarrow\quad
 w_v=c\quad\text{for every }v
\]
for some $c\in\mathbb R$. Each tested sum is finite, so no boundedness
condition on the family is needed. In particular, a subset of the places
satisfies the restricted product formula exactly when it is empty or full.
\end{theorem}
\begin{proof}
Suppose first that $K$ is a number field. For a unit of its ring of
integers all finite logarithmic terms vanish. Dirichlet's unit theorem
says that the normalized archimedean logarithm vectors span the hyperplane
$\sum_{v\mid\infty}x_v=0$; see \citet[Theorem 5.9]{MilneANT}.
The weight vector annihilates this hyperplane, hence every archimedean
weight has one value $c$.

Fix a finite prime ideal $\mathfrak p$. Finiteness of the ideal class group
\citep[Theorem 4.4]{MilneANT} gives $m>0$ and $a\in K^\times$ with
$(a)=\mathfrak p^m$. All finite logarithmic terms of $a$ except the one at
$\mathfrak p$ vanish. Subtracting $c$ times the ordinary product formula
from the weighted identity gives
\[
 (w_{\mathfrak p}-c)L_{\mathfrak p}(a)=0.
\]
Since $L_{\mathfrak p}(a)=-m\log N(\mathfrak p)\ne0$, the remaining
weight is also $c$.

Now let $K$ be a global function field, and let $C/\F_q$ be its smooth
projective curve, of genus $g$. We first recall directly why its
degree-zero divisor class group is finite. Choose an effective divisor
$A$ of degree $N\ge g$. For every degree-zero divisor $D$, Riemann--Roch
\citep[Tag 0B5B, Section 53.5]{StacksRR} gives
\[
 \dim_{\F_q}H^0(C,\mathcal O_C(D+A))\ge N+1-g\ge1.
\]
Consequently $D$ is linearly equivalent to $E-A$ for an effective divisor
$E$ of degree $N$. Only finitely many such $E$ exist: their closed points
have degree at most $N$, and the points of each degree are finite because
$C(\F_{q^d})$ is finite. This proves the required class-group finiteness.

For any two places $v,w$, the divisor
$\deg(w)[v]-\deg(v)[w]$ has degree zero. A positive multiple is principal,
so choose $a\in K^\times$ and $m>0$ with
\[
 \operatorname{div}(a)
 =m\bigl(\deg(w)[v]-\deg(v)[w]\bigr).
\]
Its weighted logarithmic relation is
\[
 m\deg(v)\deg(w)\log q\,(w_w-w_v)=0.
\]
Thus $w_v=w_w$ for every pair. Conversely constant weights satisfy the
relation by the ordinary product formula. An indicator family is constant
exactly for the empty and full subsets, proving the final assertion.
\end{proof}

A weighted logarithmic height changes under scalar multiplication by the
weighted principal sum in Theorem~\ref{thm:global-weights}. Hence, on this
standard place system, it descends to projective space exactly for
constant weights. If one adjoins a separate measured family of trivial
absolute values, their logarithmic terms vanish identically and their
weights are unrestricted. The latter construction does not alter the
rigidity assertion for the original places.


We make no claim here to have reconstructed the larger proposed arithmetic-curve or intersection-theory packages solely from the pair $(S,a)$. The semiring and its idelic unit group are exact constructions; an additional topos, a line-bundle theory, or an intersection pairing requires its own definitions and proof.

\subsection{Regular ideals and the class-group correction}
The distinction between the two zeros also changes which multipliers may be used to define ideal classes. Let $R$ be a domain with fraction field $K$, and put $S=R_0$.

\begin{proposition}[The old zero collapses unrestricted ideal classes]\label{prop:class-collapse}
On the ideals $J\ne\{\tau\}$ of $S$, the equivalence relation generated by
\[
 aJ=bJ'\qquad(a,b\in S\setminus\{\tau\})
\]
has one class. If instead one takes nonzero fractional $R$-ideals $I\subset K$, replaces them by $I_0=I\sqcup\{\tau\}\subset K_0$, and permits only multipliers in $K^\times$, the resulting class semigroup is the ordinary fractional ideal class semigroup of $R$. For a Dedekind domain this is its usual ideal class group.
\end{proposition}
\begin{proof}
Every ideal of $S$ distinct from $\{\tau\}$ contains an old ring element, and multiplication by $e=0_R$ then gives
\[
 eJ=\{\tau,e\}.
\]
The allowed choice $a=b=e$ therefore identifies every two such ideals. This proves the first assertion.

For the corrected construction, an $S$-subsemimodule $J$ of $K_0$ containing an element of $K^\times$ is $I\sqcup\{\tau\}$ with $I=J\cap K$. It contains $e$ and is closed under addition, multiplication by $R$, and additive inverses, the last by multiplication by $-1$. Thus $I$ is a nonzero $R$-submodule of $K$. The usual fractional-ideal condition, namely $dI\subset R$ for some $d\in R\setminus\{0\}$, is exactly $dJ\subset S$. Conversely every such $I$ gives an $S$-subsemimodule. Products are the finite sums of pairwise products in both descriptions, so $(IJ)_0=I_0J_0$. Multiplication by $a\in K^\times$ also commutes with this correspondence. Consequently the two class equivalence relations and their products agree. For a Dedekind domain the nonzero fractional ideals are invertible, giving the asserted group; see \citet[Theorem 3.20, p.\ 53]{MilneANT}.
\end{proof}

For a concrete failure of the unrestricted definition, take $R=\Z[\sqrt{-5}]$. The ideal $(2,1+\sqrt{-5})$ has index two: reduction modulo it maps $\sqrt{-5}$ to $1$ in $\F_2$. It is not principal, since a generator would have norm two, but $a^2+5b^2=2$ has no integer solution. Thus its ordinary class is nontrivial, whereas the first equivalence in Proposition~\ref{prop:class-collapse} identifies it with the unit ideal. If regular multipliers are imposed while the old-zero ideal $\{\tau,e\}$ is retained, that ideal contributes an additional absorbing class; restricting the ideals as well is necessary.

\begin{proposition}[Valuation subsemirings]\label{prop:valuations}
Let $V\subset K_0$ be a unital subsemiring containing $\tau$ and $e$. Suppose that for every $a\in K^\times$, at least one of $a,a^{-1}$ belongs to $V$. Then $V=W_0$ for a valuation ring $W$ of $K$. Conversely every valuation ring gives such a subsemiring. For $K=\Q$, these are $\Q_0$ and $(\Z_{(p)})_0$ for rational primes $p$.
\end{proposition}
\begin{proof}
Applying the hypothesis to $-1$ shows that $-1\in V$. Therefore $W=V\cap K$ is a ring and the stated inverse condition is exactly the defining valuation-ring condition. The converse follows from the same definition.

For the rational classification, every subring $W\subset\Q$ containing $\Z$ is a localization of $\Z$ at a set of primes. Indeed, if a reduced fraction $a/b$ lies in $W$, Bezout's identity gives $1/b\in W$, hence the inverses of every prime dividing $b$ lie in $W$. If two distinct primes $p,q$ were not inverted, neither $p/q$ nor $q/p$ would belong to $W$, contradicting the valuation condition. At most one prime is therefore not inverted, giving exactly the rings asserted.
\end{proof}

Finally, the full-support copy of $\A_\Q$ inside $\mathcal A_0$ preserves addition, multiplication, and the multiplicative identity, but sends its additive zero to the old, fully supported zero. It is not a zero-preserving homomorphism of semirings. The zero-adjoined ring $(\A_\Q)_0$ does map into $\mathcal A_0$, with its new zero sent to empty support. This distinction is required whenever the adelic inclusion is used in a categorical assertion.


\section{Linear symmetry and projective coordinate counts}\label{sec:linear}
These elementary consequences specify another part of the geometry that is intrinsic to zero adjunction.

\begin{proposition}\label{prop:gl}
For every nonzero or zero commutative ring $R$ and $n\ge1$,
\[
 \operatorname{GL}_n(R_0)\cong(R^\times)^n\rtimes\mathfrak S_n.
\]
Every invertible matrix is a permutation matrix with one old ring unit in each occupied position and $\tau$ elsewhere.
\end{proposition}
\begin{proof}
Apply the support character $R_0\to\B$ to an inverse pair of matrices. Let their Boolean matrices be $A,B$, with $AB=BA=I$. If $a_{ij}=1$, row $j$ of $B$ can have a $1$ only in column $i$, since off-diagonal entries of $AB$ are zero. It cannot be empty, since $BA=I$. If row $i$ of $A$ contained two ones, the corresponding two rows of $B$ would both be supported in column $i$, and $BA$ would have an off-diagonal one. Thus every row of $A$ has one one, and similarly every column does. The supported entries of the original matrix and its inverse must then multiply to $1_R$. Conversely a permutation of unit entries is plainly invertible.
\end{proof}

\begin{proposition}[A coordinate-orbit counting functor]\label{prop:points}
For a finite field $\F_q$, take vectors in $(\F_q)_0^{n+1}$ whose entries generate the unit ideal, and identify vectors differing by a common unit of $\F_q$. Their number is
\[
 C_n(q)=\frac{(q+1)^{n+1}-2^{n+1}}{q-1}
       =\sum_{j=0}^{n}a_{n,j}q^j,
 \qquad a_{n,j}=\sum_{k=j+1}^{n+1}\binom{n+1}{k}.
\]
For $n=1$ the count is $q+3$. Its associated formal Tate Euler product is $\prod_{j=0}^{n}\zeta(s-j)^{a_{n,j}}$.
\end{proposition}
\begin{proof}
Fix the nonempty support $I$ of the vector, of cardinality $k$. Its amplitudes form a vector of $\F_q^k$. The unit-ideal condition is exactly that this amplitude vector is not zero. Multiplication by $\F_q^\times$ acts freely on those $q^k-1$ vectors. Thus this support contributes $(q^k-1)/(q-1)=1+q+\cdots+q^{k-1}$. Sum over the $\binom{n+1}{k}$ supports. Summing the geometric series in the other order gives the coefficients $a_{n,j}$. The Euler-product assertion is the determinant calculation for these Tate multiplicities, not an identification with a different scheme's point-counting functor.
\end{proof}

For example, in two coordinates $[1:\tau]$ and $[1:e]$ are distinct. The first has one coordinate present; the second has two. The stated functor is deliberately specified by its vectors and equivalence relation. A comparison with a particular blueprint projective scheme requires matching that scheme's notion of points; a numerical equality alone is not that comparison.

\section{Lattice support, projective coordinates, and cycle invariants}
\label{sec:lattice-coordinate-geometry}
Let $L$ be a nontrivial bounded distributive lattice and $K$ a field. Write
\[
 G_L(K)=\{(0,\lambda):\lambda\in L\}\cup\{(a,1):a\in K\}
 \subseteq K\times L,
\]
where addition uses join in the second coordinate and multiplication uses
meet. The repeated element $(0,1)$ is included only once. These operations
are closed: a nonzero amplitude can occur only at top support. The new
zero is $(0,0)$, and the unit is $(1,1)$. Thus this is a semiring, with
amplitude and support homomorphisms to $K$ and $L$. When $L=\B$ it is
$K_0$. We call a vector unimodular when its entries generate the unit ideal,
and take projective coordinates to mean its orbit under common scalar
units. This specifies an orbit functor without asserting a projective
scheme realization.

\begin{proposition}[Support profiles]\label{prop:lattice-profiles}
A vector of length $N$ over $G_L(K)$ is unimodular exactly when at least
one amplitude is nonzero. Its scalar orbit is specified by a nonempty
set $A\subseteq\{1,\ldots,N\}$ of top-support positions, a point of
$\mathbf P^{|A|-1}(K)$, and an arbitrary proper lattice element at every
position outside $A$. If $|L|=\ell<\infty$ and $K=\F_q$, the orbit count is
\[
 C_{N,L}(q)=\frac{(q+\ell-1)^N-\ell^N}{q-1}
 =\sum_{j=0}^{N-1}a_jq^j,
 \quad
 a_j=\sum_{r=j+1}^{N}\binom Nr(\ell-1)^{N-r}.
\]
Its formal point-counting zeta series is
\[
 \exp\!\left(\sum_{m\ge1}C_{N,L}(q^m)\frac{T^m}{m}\right)
 =\prod_{j=0}^{N-1}(1-q^jT)^{-a_j}.
\]
\end{proposition}
\begin{proof}
A coordinate of nonzero amplitude is a unit, so suffices for unimodularity.
If all amplitudes vanish, every linear combination has amplitude zero
and cannot be the unit. Fixing the top-support set leaves a nonzero vector
of $K^A$, on which $K^\times$ acts by the usual scalar action, and leaves
the proper lattice decorations fixed. The action on that nonzero vector
is free. A set $A$ of size $r$ therefore contributes
$(\ell-1)^{N-r}(q^r-1)/(q-1)$. Sum over $A$, and expand the geometric
sum to obtain the coefficients $a_j$. The last identity follows from
$\sum_{m\ge1}z^m/m=-\log(1-z)$ as a formal power series.
\end{proof}

\begin{proposition}[Linear symmetry and cycles]\label{prop:lattice-cycles}
For every nonzero commutative ring $R$,
\[
 \operatorname{GL}_n(G_L(R))=(R^\times)^n\rtimes\mathfrak S_n.
\]
For a matrix $A$ over $G_L(R)$, the signed Leibniz determinant has ordinary
determinant amplitude and support equal to the lattice permanent of its
support matrix. The $k$ th characteristic coefficient has support the join
of the support products of all principal cycle covers on $k$ vertices.
For a word in several matrices, its trace has support the join of the
corresponding colored closed-walk products. All these expressions are
invariant under monomial simultaneous conjugation.

Over a field, conjugacy classes of invertible monomial matrices under monomial conjugation are classified by the multiset
of permutation-cycle lengths and cycle products. A cycle of length $r$
and product $w$ contributes $t^r-w$ to the characteristic polynomial and
contributes $r w^{m/r}$ to the trace of the $m$ th power when $r\mid m$,
and zero otherwise. Block replication $A\mapsto A\otimes I_d$ raises
the split determinant and characteristic polynomial to the $d$ th power
and multiplies trace powers by the old scalar $d$.
\end{proposition}
\begin{proof}
Let $AB=BA=I$. In each diagonal amplitude sum some term is nonzero,
since the sum is $1$. Its two factors, say $A_{ik},B_{ki}$, have top support.
The off-diagonal support equations in $AB$ force $B_{kj}=0$ for $j\ne i$;
those in $BA$ force $A_{i l}=0$ for $l\ne k$. Here $0$ means the new zero.
The remaining diagonal equation says $A_{ik}B_{ki}=1$. Distinct rows have
distinct such columns, proving the monomial description. The converse is
immediate. If $R$ is the zero ring, this argument does not apply: a product
Boolean lattice can have different permutations on its different factors.

Apply amplitude and support to the Leibniz expansion. Signs are old units,
so have top support. This gives the determinant and principal-minor
formulas. Expanding a trace word lists exactly its colored closed walks.
A diagonal change of frame supplies a ratio of vertex units to each edge;
these ratios cancel on every cycle. A permutation relabels vertices,
proving invariance.

For conjugacy, a permutation change of basis identifies cycles of equal
length. Along a cycle choose one vertex scalar and successively determine
the others to match edge weights. The last edge agrees exactly when the
two cycle products agree. Ordering a basis by cycles gives the stated
characteristic and trace formulas directly. Finally a permutation of the
tensor-product basis turns $A\otimes I_d$ into $d$ diagonal copies of $A$,
with the new zero off the blocks. Matchings cannot cross those blocks,
so determinants and characteristic polynomials multiply, while traces add.
\end{proof}

A lattice join need not be an attained matching level. In $L=\B^2$, give
the diagonal entries of a $2$ by $2$ zero-amplitude matrix support $(1,0)$
and the off-diagonal entries support $(0,1)$. The two matching levels join
to $1$, although neither matching has level $1$. For Boolean support the
join does express existence of a matching; for a general lattice different
Boolean branches may have different witnesses.

\section{Measured adelic supports and logarithmic descent}
\label{sec:measured-defects}
A measured adelic family consists of a field $K$, a measure space
$(\Omega,\mathcal A,\nu)$, and absolute values $|\cdot|_\omega$ such that
$\ell_a(\omega)=\log|a|_\omega$ belongs to $L^1$ for every $a\in K^\times$.
It is proper when every $\ell_a$ has integral zero. This is the
Chen--Moriwaki framework \citep[Section 3.1]{ChenMoriwakiAdelic}. Let
$\mathcal B=\mathcal A/\text{null}$ be its Boolean measure algebra.
Define
\[
 \Delta([A])(a)=\int_A\ell_a\,d\nu,\qquad
 \mathcal D(f)(a)=\int_\Omega f\ell_a\,d\nu
 \quad(f\in L^\infty).
\]
Call $[A]$ balanced when $\Delta([A])=0$.

\begin{theorem}[Defect transform and localization]\label{thm:measured-defect}
The map $\Delta$ is a valuation from $\mathcal B$ to the additive group
$\operatorname{Hom}(K^\times,\mathbb R)$:
\[
 \Delta(A\cup B)+\Delta(A\cap B)=\Delta(A)+\Delta(B).
\]
The map $\mathcal D$ is real-linear, and $\mathcal D(1_A)=\Delta(A)$.
For a proper family, constants lie in its kernel and
$\Delta(A^c)=-\Delta(A)$. The bilinear pairing
\[
 \left(L^\infty/\ker\mathcal D\right)
 \times\operatorname{span}_{\mathbb R}\{\ell_a:a\in K^\times\}
 \longrightarrow\mathbb R,
 \qquad([f],g)\longmapsto\int fg\,d\nu
\]
is nondegenerate in both variables.

There is a canonical localization isomorphism
\[
 G_{\mathcal B}(K)[(0,[A])^{-1}]
 \cong\downarrow[A]\cong\mathcal B_{\Omega|A}.
\]
The restricted adelic family is proper exactly when $[A]$ is balanced.
A measure-preserving refinement preserving the absolute values commutes
with these defect maps and localization identifications.
\end{theorem}
\begin{proof}
Null changes do not affect integrals, and $\ell_{ab}=\ell_a+\ell_b$ proves
the character assertion. The indicator identity for union and intersection
proves the valuation formula. Bounded functions multiply integrable
functions into $L^1$, so $\mathcal D$ is well defined and linear. Properness
and $1_{A^c}=1-1_A$ give the remaining formulas.

The pairing is well defined by the definition of the kernel. A class
pairing trivially with every $\ell_a$ is zero. Conversely a nonzero
$g\in L^1$ pairs positively with the bounded function $\operatorname{sgn}g$,
since $\int|g|>0$. This proves nondegeneracy without claiming identification
with a full continuous or algebraic dual.

For any semiring and multiplicative idempotent $z$, localization at $z$
is $zS$, with unit $z$: every map inverting $z$ sends it to $1$ and factors
uniquely through $s\mapsto zs$. Here multiplication by $(0,[A])$ erases
amplitude and intersects support with $[A]$, giving $\downarrow[A]$.
Measurable subsets of $A$, modulo their restricted null sets, are exactly
that relative Boolean algebra. Its product formula is precisely
$\Delta(A)=0$. For a refinement $\pi$, change of variables gives
$\mathcal D'(f\circ\pi)=\mathcal D(f)$ and inverse images commute with
intersection, proving both naturality statements.
\end{proof}

The quotient in this theorem is a vector space. If the family is proper
and $\mathcal D\ne0$, its kernel contains $1$ but is proper, so it is not an
ideal of $L^\infty$. No quotient algebra or induced Boolean idempotent
structure is asserted for that vector-space quotient.

\begin{proposition}[Heights, neutral places, and metric states]
\label{prop:height-neutrality}
For $f\in L^\infty$ and $\mathbf x\in K^{n+1}\setminus\{0\}$, put
\[
 h_f(\mathbf x)=\int f(\omega)\log\max_i|x_i|_\omega\,d\nu.
\]
This is finite, and
$h_f(a\mathbf x)-h_f(\mathbf x)=\mathcal D(f)(a)$.
Consequently it descends to $\mathbf P^n(K)$ exactly when
$f\in\ker\mathcal D$. The same assertion, with the opposite sign,
holds for $-\int f\log\|s\|_\omega\,d\nu$ whenever the norm logarithms
of nonzero global vectors are integrable.

Adjoining any measured family $T$ of trivial absolute values leaves these
heights and canonical trivial-norm degrees unchanged. It replaces
$\mathcal B$ by $\mathcal B\times\mathcal B_T$ and the defect kernel by
$\ker\mathcal D\oplus_\infty L^\infty(T)$. In particular, for the standard
global place system, balanced supports after this extension are exactly
an empty or full original place set together with an arbitrary added
support.

Up to positive rescaling, the local states consisting of absence, a zero
seminorm on a line, and a nondegenerate seminorm form the chain
$0<1<2$, whose tensor product is meet. Measurable such states modulo null
sets form the bounded distributive lattice
\[
 \{(P,N)\in\mathcal B^2:N\le P\},
\]
where $P$ records presence and $N$ nondegeneracy. The function
$(P,N)\mapsto\Delta(N)$ is a valuation on this lattice.
\end{proposition}
\begin{proof}
Discard zero coordinates. The logarithm of the maximum is the maximum of
finitely many $L^1$ functions, and its absolute value is bounded by the sum
of their absolute values. Scalar multiplication adds $\ell_a$ to that
maximum. The norm formula follows in exactly the same way from
$\|as\|_\omega=|a|_\omega\|s\|_\omega$. Invariance under every scalar
is equivalent to vanishing of the displayed character, including when
$n=0$.

On the added family every principal logarithm is zero and the maximum
of the absolute values of a nonzero tuple is $1$. The canonical trivial
norm is $1$ on nonzero vectors. Thus all additional integrands vanish.
Measurable sets and essentially bounded functions on a disjoint union
split into their two components, yielding the product and kernel formulas.
The global-field assertion follows from Theorem~\ref{thm:global-weights}.

A seminorm on a real line is either zero or a positive multiple of any
fixed norm. Tensoring with absence remains absent, tensoring a present
zero seminorm gives a zero seminorm, and tensoring two nondegenerate
seminorms remains nondegenerate. These are the three meet rules. The
threshold maps $j\mapsto1_{j\ge1}$ and $j\mapsto1_{j\ge2}$ identify the
chain with the incidence pairs in $\B^2$. Apply this pointwise; the inverse
assigns $2$ on $N$, $1$ on $P\setminus N$, and $0$ outside $P$. Projection to $N$
preserves joins and meets, so the valuation identity follows. This
incidence lattice is generally not Boolean.
\end{proof}

Booleanization of an absolute value itself only tests whether a scalar
vanishes: $|a|_\omega>0$ for every $a\ne0$. It is therefore independent of
the place. A trivial absolute value is nondegenerate and differs from an
identically zero seminorm.

The measure algebra above is always Boolean sigma-complete, but need not
be complete. To see this, equip an uncountable set $X$ with its
countable/cocountable sigma-algebra and counting measure, and use the
trivial absolute value at every point. All principal logarithms vanish,
so the family is proper. Its null ideal is $\{\varnothing\}$. If $Y$ and
$X\setminus Y$ are both uncountable, the singleton subsets of $Y$ have no
least measurable upper bound: any such upper bound is cocountable and
can be made strictly smaller by deleting a point outside $Y$. Thus the
unqualified completeness claim fails. The displayed constructions use
only finite lattice operations and require no completeness assumption.

\section{Monomial bundles and support-refined invariant gluing}
\label{sec:monomial-invariant-gluing}
\begin{proposition}\label{prop:cover-local-system}
In the usual local-system category of a connected space, rank $n$ local
systems with monomial transition maps are equivalent to degree $n$ finite
covering local systems equipped with a rank-one local system on the cover.
For coefficients $\mathbb C_0$ these are exactly the local systems defined
by invertible $n$ by $n$ split matrices. Their ordinary loop determinants
factor over the cycles of the covering permutation.
\end{proposition}
\begin{proof}
The permutation parts transport the $n$ basis labels and define the cover.
On each label the diagonal multiplier transports a one-dimensional fiber,
defining the line local system. Conversely the direct sum over the sheets
has monomial transition maps. The constructions are inverse on objects
and on their compatible isomorphisms. The split assertion is the general
linear-group calculation, and the loop determinant is the cycle formula
in Proposition~\ref{prop:lattice-cycles}.
\end{proof}

For a sheaf of modules or line bundle $V$, write $\operatorname{Spl}(V)$ for
the sheaf coproduct of $V$ and an absent point. A section is locally either
an old section or absent; addition uses the old sum when both are present,
and absence is its identity. Products take values in the corresponding
tensor product and are absent if either factor is absent.

\begin{proposition}[Matrix Higgs data]\label{prop:matrix-higgs-support}
Let $X$ be a smooth complex curve. Consider a vector bundle presented by
monomial transition matrices, together with local $n$ by $n$ matrices of
sections of $\operatorname{Spl}(K_X)$ that transform by conjugation.
The split characteristic coefficients define sections
\[
 c_k\in\Gamma(X,\operatorname{Spl}(K_X^{\otimes k})),
 \qquad1\le k\le n.
\]
Their amplitudes are the ordinary characteristic coefficients. At a
geometric point $c_k$ is absent exactly when no principal $k$-vertex cycle
cover is present; it is the old zero exactly when some such cover is
present and the amplitude coefficient vanishes.
\end{proposition}
\begin{proof}
Expand $c_k$ as the signed sum of its principal cycle-cover products. Each
product has tensor degree $k$. A diagonal change of bundle frame cancels
around each cycle; a permutation relabels the covers. Thus the local
expressions agree on overlaps in $K_X^{\otimes k}$ and glue in its split
sheaf. Applying amplitude and Boolean support to the expansion gives the
two assertions at a geometric point.
\end{proof}

These are arbitrary matrix Higgs fields on a monomially reduced bundle.
They are not restricted to the Lie algebra of the monomial normalizer,
which would impose diagonal fields. A further restriction of a geometric
Langlands functor requires a specified spectral stack, functor and all of
its hypotheses; the invariant-gluing proposition supplies none of those
additional structures.

\begin{proposition}[A split torsor calculation]\label{prop:split-torsor}
Let $M$ be a nonempty simply transitive $W$-set, and adjoin two distinct
$W$-fixed points. Its action groupoid is equivalent to
$BW\sqcup BW\sqcup *$. Finite products give the corresponding product
groupoids and a Boolean support map recording absence at each factor.
\end{proposition}
\begin{proof}
The two fixed orbits each have stabilizer $W$ and give $BW$. A simply
transitive orbit has trivial stabilizer, so its action groupoid is
contractible. The decomposition is invariant under the action, and
products of these decompositions prove the final assertion. Applying this
calculation to a proposed arithmetic torsor still requires that torsor to
have been defined and its simple transitivity proved.
\end{proof}


\subsection{Support reductions and the information they retain}

Several support constructions require only elementary algebra. Stating
their precise targets prevents them from being read as reconstructions
of the original arithmetic data.

\begin{proposition}[A multiplicative support quotient]
\label{prop:rees-support-quotient}
For a unital monoid $M$, the Rees quotient of $M\times\B$ by its absorbing
ideal $M\times\{0\}$ is $M$ with one new multiplicative absorber.
\end{proposition}
\begin{proof}
Collapse all $(m,0)$ to a single point $*$ and retain the classes of
$(m,1)$ separately. The product of two retained classes is $(mn,1)$;
any product involving the collapsed ideal is in that ideal. These are
exactly the asserted operations, and $(1,1)$ is the unit. This is a
statement about multiplicative monoids. No unspecified additive
operation or semiring isomorphism follows from it.
\end{proof}

\begin{proposition}[Initial data do not reconstruct a valued element]
\label{prop:initial-data}
Let $K$ be a valued field, with associated graded ring
$\operatorname{gr}_vK$. On $K_0$, the new zero and the old zero remain
distinct. For an old nonzero $x$, its valuation and initial class
$\operatorname{in}_v(x)\in(\operatorname{gr}_vK)_{v(x)}$ specify its first
graded term, but need not determine $x$. If a Laurent-series coordinate
is fixed, the initial class is equivalently the leading coefficient
together with its exponent.
\end{proposition}
\begin{proof}
For $x\ne0$, the initial class is its image in
$\{y:v(y)\ge v(x)\}/\{y:v(y)>v(x)\}$. Thus two elements of the same
valuation have the same initial class whenever their difference has
larger valuation. In $k((t))$, the distinct elements $t$ and $t+t^2$
have valuation $1$ and initial class $t$. They also have the same
occupied support in $K_0$. The two zero states are separately specified
by construction. Consequently a datum consisting of support, valuation,
and leading coefficient is initial data, not a complete invariant of
a Laurent series or any larger valued-field realization containing
this example.
\end{proof}

For a finite witness set $W$, its cardinality can be encoded inside
$\mathbb Z_0$ by the sum $c(W)=\sum_{w\in W}1$: this is the new zero for
$W=\varnothing$ and the old positive integer $|W|$ otherwise. Hence its
Boolean support is exactly the truth value of $W\ne\varnothing$.
Disjoint unions and Cartesian products give addition and multiplication
of these counts, including empty sets. Applied to any explicitly
specified finite Diophantine witness set, this records solvability
exactly, but does not by itself prove that a witness exists. It uses the
natural copy of $\mathbb N$ in $\mathbb Z_0$ whose zero is the new zero;
it makes no assertion that the larger zero-adjoined set $\mathbb N_0$
is an inverse semiring.


\section{Associative effective coefficients}\label{sec:criterion}
Let $B$ be a finite named basis containing a unit $1$. Suppose integers $c_{ij}^k(q)\in\Z[q,q^{-1}]$ give an associative unital algebra on $B$.

\begin{proposition}[The dimension criterion]\label{prop:criterion}
Effective coefficients $a_{ij}^{k,+}+\eps a_{ij}^{k,-}$ with signed traces $c_{ij}^k$ define an associative unital algebra over $\mathcal E$ if and only if their dimensions $d_{ij}^k=a_{ij}^{k,+}+a_{ij}^{k,-}$ define an associative unital algebra and satisfy, coefficientwise,
\[
 d_{ij}^k\ge |c_{ij}^k|,\qquad d_{ij}^k\equiv c_{ij}^k\pmod2.
\]
The unit coefficients are required to be the ordinary Kronecker coefficients in both tables. In that case the lift is unique for the given pair of tables and is
\[
 a_{ij}^{k,\pm}=\frac{d_{ij}^k\pm c_{ij}^k}{2}.
\]
\end{proposition}
\begin{proof}
Evaluation at $\eps=1$ and $\eps=-1$ embeds
\[
 \Z[q,q^{-1},\eps]/(\eps^2-1)
 \longrightarrow\Z[q,q^{-1}]\times\Z[q,q^{-1}],
 \qquad a^++\eps a^-\longmapsto(d,c).
\]
It is injective because the inverse formulas are $(d+c)/2$ and $(d-c)/2$; its image consists of pairs congruent modulo two. The associativity equations and the unit equations are polynomial equations in the coefficients. They hold in the left-hand ring if and only if they hold after both evaluations. Nonnegativity of the inverse formulas is equivalent to the displayed inequalities. This proves both directions.
\end{proof}

This proposition is an algebraic lifting criterion. If it is realized inside a category of matrices of graded sign representations, the natural Grothendieck group is the \emph{split} Grothendieck group of the additive monoidal subcategory generated by those matrices. That construction does not automatically give an abelian, rigid, or idempotent-complete category with the basis objects simple. We use no such stronger categorification claim.

\subsection{Character values with a prescribed support}\label{sec:character-support}
The dimension--trace pair used above is the two-character case of a standard finite-group calculation. Let $\Gamma$ be finite, let $R(\Gamma)$ be its complex representation ring, and write $R(\Gamma)_0$ for its zero adjunction. Character evaluation at $g\in\Gamma$ is a ring homomorphism and extends to
\[
 \widehat{\operatorname{ev}}_g:R(\Gamma)_0\longrightarrow\mathbb C_0,
 \qquad \tau\longmapsto\tau.
\]
Every present representation class maps to a present character value, even when that value is zero.

\begin{proposition}[Characters on a prescribed set of grades]\label{prop:character-support}
Let $E$ be an index set and $B\subseteq E$ a chosen set of present grades. For each $b\in B$, let $f_b$ be a complex class function on $\Gamma$. There is a unique isomorphism class of finite-dimensional genuine representation $V_b$ with character $f_b$ at each present grade if and only if
\[
 m_{b,\chi}=\frac1{|\Gamma|}\sum_{g\in\Gamma}f_b(g)\chi(g^{-1})
 \in\mathbb Z_{\ge0}
 \qquad(\chi\in\operatorname{Irr}(\Gamma)).
\]
The resulting zero-adjoined coefficient is $[V_b]$ on $B$ and $\tau$ on $E\setminus B$. A zero character on a present grade gives the old zero class, not $\tau$. The full family of zero-adjoined character evaluations is injective on $R(\Gamma)_0$.
\end{proposition}
\begin{proof}
The irreducible characters form an orthonormal basis of the class functions, and the displayed inner product computes each multiplicity \citep[\S2.3, Theorem~4 and Corollary~2, p.~16; \S2.5, Theorem~6, p.~19]{SerreCharacters}. Thus nonnegative integral values construct
\[
 V_b=\bigoplus_{\chi\in\operatorname{Irr}(\Gamma)}
 V_\chi^{\oplus m_{b,\chi}},
\]
and character uniqueness proves necessity and uniqueness. All multiplicities vanish exactly for the zero representation. The chosen support convention places this class at a present grade, while absent grades have coefficient $\tau$. Finally, ordinary characters distinguish elements of $R(\Gamma)$, and every extended evaluation distinguishes the new zero from a present class. Together these facts prove injectivity.
\end{proof}

Replacing $\chi(g^{-1})$ by $\chi(g)$ in the formula extracts the multiplicity of the dual irreducible. A constant boundary class function equal to $-2$ is $-2$ times the trivial character and cannot represent a genuine module, since its value at the identity is negative. These are conditional coefficient tests. A prescribed zero grade or a virtual principal part does not supply a new existence theorem for a moonshine module, and the chosen set $B$ cannot be recovered from an ordinary character series after its zero coordinates have been collapsed.


\section{The rank-three matrix algebra}\label{sec:rank3}
This section defines the algebra and computes its multiplication without using a precomputed table. Work first over $K=\Q(D)$. For $\lambda\in K$ set
\[
 P_1(\lambda)=\begin{pmatrix}D&\lambda\\0&0\end{pmatrix},\qquad
 P_2(\lambda)=\begin{pmatrix}0&0\\1&D\end{pmatrix}.
\]
They satisfy $P_i^2=DP_i$ and $P_iP_jP_i=\lambda P_i$ for $i\ne j$. Let $W=\{1,2,12,21\}$ and let $P_w$ be the corresponding product. Put
\[
 \lambda_1=D,\quad\lambda_2=(D-2)^2,\quad
 \delta_i=\lambda_i-\lambda_{3-i},
\]
\[
 G_i=D(D^2-\lambda_{3-i}),\qquad
 S=D(D^2-\lambda_1)(D^2-\lambda_2)=4D^2(D-1)^2.
\]
In $K\times K\times M_2(K)\times M_2(K)$ define
\[
 1=(1,1,I,I),\qquad\Sigma=(0,S,0,0),
\]
and define $\gamma_w^i$ to have first component zero, second component $D^{|w|-1}G_i$, matrix component $\delta_iP_w(\lambda_i)$ in block $i$, and zero in the other block. These ten elements are a basis: the first scalar component isolates $1$; each set of four matrices $P_w(\lambda_i)$ is independent when $\lambda_i(D^2-\lambda_i)\ne0$; and the remaining second scalar component isolates $\Sigma$. The determinant of the four columns obtained by listing the entries of $P_1,P_2,P_{12},P_{21}$ is $\lambda(\lambda-D^2)^2$, proving the required independence.

The cell parameters and the bottom scalar character reproduce the rank-three formulas of \citet[\S5.1 and Figure 2]{Mulmuley}. His printed display (17) repeats $b_1$ for both cell superscripts; taken literally that display is not ten distinct elements. We therefore state the matrix realization and all parameters explicitly. The unimodality correction below uses a coefficient printed independently in Figure 2, and does not depend on resolving that repeated-subscript typography.

For arbitrary cell parameters $\lambda_i,\delta_i,G_i,S$ with a nonsingular basis of this form, set
\[
 W_0=\frac{G_1G_2}{S},\qquad
 U_i=\frac{G_i(G_i-\delta_iD)}{S},\qquad
 V_i=\frac{G_i(DG_i-\delta_i\lambda_i)}{S}.
\]
In each $2$ by $2$ block, reduction of a product gives
\[
 P_xP_y=D^\alpha\lambda_i^\beta P_z,
 \quad \beta\in\{0,1\},\quad
 |x|+|y|=|z|+\alpha+2\beta.
\]
All these exponents are obtained just by removing a repeated letter and then an alternating triple.

\begin{lemma}[The complete multiplication rule]\label{lem:table}
Apart from unit products, the multiplication is
\begin{align*}
 \gamma_x^i\gamma_y^i
 &=\delta_iD^\alpha\lambda_i^\beta\gamma_z^i
   +D^{|z|+\alpha-1}
      \begin{cases}U_i/D,&\beta=0,\\V_i,&\beta=1\end{cases}\Sigma,\\
 \gamma_x^1\gamma_y^2&=D^{|x|+|y|-2}W_0\Sigma,\\
 \Sigma\gamma_x^i&=\gamma_x^i\Sigma=D^{|x|-1}G_i\Sigma,
 \qquad \Sigma^2=S\Sigma.
\end{align*}
The cross-cell rule also holds with the factors reversed.
\end{lemma}
\begin{proof}
The matrix blocks of a within-cell product determine its cell coefficient as $\delta_iD^\alpha\lambda_i^\beta$. Subtracting that multiple of $\gamma_z^i$ leaves only the second scalar component. Its coefficient at $\Sigma$ is
\[
 \frac{D^{|x|+|y|-2}G_i^2
       -\delta_iD^{\alpha+|z|-1}\lambda_i^\beta G_i}{S}.
\]
Use the length identity to factor $D^{|z|+\alpha-1}$. If $\beta=0$, the remaining factor is $U_i/D$; if $\beta=1$, it is $V_i$. Cross-cell products have both matrix blocks zero, so their scalar coefficient is $D^{|x|+|y|-2}G_1G_2/S$. Multiplication by $\Sigma$ is read from the second scalar component. These computations account for every component and hence every coefficient.
\end{proof}

For the trace parameters, the simplified bottom factors are
\[
 W_0=U_1=U_2=D,\quad V_1=5D-4,\quad V_2=2D^2-5D+4.
\]
In particular, all bottom coefficients have nonnegative Laurent coefficients after $D=q+2+q^{-1}$. The potential signs occur in $\delta_i$.

\begin{proposition}[The unimodality correction]\label{prop:unimodality}
The coefficient of $\gamma_1^2$ in $\gamma_{12}^2\gamma_1^2$ is
\[
 c=(D-1)(D-4)(D-2)^2.
\]
Writing $t=q^{1/2}-q^{-1/2}$, its positive bar-invariant factor after removing its vanishing factor is
\[
 c=t^2(q+1+q^{-1})(q+q^{-1})^2.
\]
The factor has coefficients $(1,1,3,2,3,1,1)$ and is not unimodal. There are six within-cell coefficients with this same value in the multiplication rule.
\end{proposition}
\begin{proof}
The identity $P_1P_2P_1=\lambda_2P_1$ gives the coefficient $\delta_2\lambda_2$, which is the displayed $c$. Substitute $D-4=t^2$, $D-1=q+1+q^{-1}$, and $D-2=q+q^{-1}$, and multiply the two elementary Laurent polynomials. The sequence increases from $1$ to $3$ and then decreases to $2$ before increasing again, so it is not unimodal.

For completeness the six ordered word pairs with $\alpha=0$, $\beta=1$ are $(1,21)$, $(2,12)$, $(12,1)$, $(21,2)$, $(12,12)$, and $(21,21)$. Their coefficients are exactly $\delta_2\lambda_2$. If $c=\pm t^a f$ with $f$ nonzero and having nonnegative coefficients, then $f(1)>0$. The order of vanishing at $q=1$ forces $a=2$. Thus changing the extracted power cannot repair unimodality of this positive factor. Figure 2 of \citet{Mulmuley} prints this same rational expression in $q$, so the correction concerns the displayed normalization there.
\end{proof}

\section{Effective associative lifts}\label{sec:lifts}
This section proves Theorem~\ref{thm:lifts-main}. Keep $D=q+2+q^{-1}$. In the dimension table put
\[
 \lambda_1^+=D(4D-9),\qquad\lambda_2^+=3D^2-10D+4,
\]
and choose either
\[
 \delta^+=D(D-3)\quad\hbox{or}\quad\delta^+=(D-1)(D-2).
\]
Use this same $\delta^+$ in both matrix blocks, and set
\[
 G_1^+=10D\delta^+,\quad G_2^+=15D\delta^+,\quad
 S^+=30D(\delta^+)^2.
\]
These give
\[
 W_0^+=5D,\quad U_1^+=3D,\quad U_2^+=7D,
\]
\[
 V_1^+=D(2D+3),\qquad V_2^+=6D^2+5D-2.
\]
The basis is nonsingular over $\Q(D)$ and Lemma~\ref{lem:table} shows that every structure coefficient is an integral Laurent polynomial. It is associative because it is multiplication in a product of matrix algebras.

\begin{lemma}\label{lem:positive}
Each of the two dimension tables has nonnegative coefficients, dominates the coefficientwise absolute value of the trace table, and is congruent to it modulo two.
\end{lemma}
\begin{proof}
Write $\delta=(D-1)(D-4)$. Its Laurent coefficient sequence from weight $-2$ to $2$ is $(1,-1,0,-1,1)$. The first $\delta^+$ is its coefficientwise absolute value $(1,1,0,1,1)$; the second is $(1,1,2,1,1)$. Both dominate $|\delta|$ and are congruent to both $\delta$ and $-\delta$ modulo two.

The dimension sequences of $\lambda_1^+$ and $\lambda_2^+$ are $(4,7,6,7,4)$ and $(3,2,2,2,3)$. They dominate those of $\lambda_1=D$ and $\lambda_2=(D-2)^2$, and the differences have even coefficients. The elementary inequality $|ab|\le |a|\,|b|$ coefficientwise for Laurent polynomials proves the within-cell assertion, using Lemma~\ref{lem:table}.

For bottom coefficients it suffices to check the factors $W_0,U_i,V_i,G_i,S$, since the remaining factor is a nonnegative power of $D$. The differences for $W_0,U_1,U_2$ are $4D,2D,6D$. The differences for $V_1,V_2$ are $2D^2-2D+4$ and $4D^2+10D-6$, each with nonnegative even Laurent coefficients. For the first choice of $\delta^+$, direct multiplication gives the following difference sequences in increasing weight order:
\[
 G_1^+-G_1=(10,26,18,4,18,26,10),
\]
\[
 G_2^+-G_2=(14,40,34,16,34,40,14).
\]
Both ranges are $[-3,3]$. These are the expansions of the explicit polynomials
\[
 10D^2(D-3)-4D(D-1),\qquad
 15D^2(D-3)-D^2(D-1).
\]
The difference $S^+-S$ has sequence
\[
 (30,116,156,112,150,216,150,112,156,116,30)
\]
on $[-5,5]$. This completes the finite coefficient check. Congruence modulo two follows directly from $G_1^+\equiv0\equiv G_1$, $G_2^+\equiv D^2(D-3)\equiv D^2(D-1)$, and $S^+\equiv0\equiv S$. For the second choice, $\delta^+$ is increased by $2$, so the corresponding $G_i$ and $S$ increase coefficientwise by even nonnegative Laurent polynomials. This proves the remaining case.
\end{proof}

\begin{proof}[Construction and verification of the first two lifts]
Associativity of both tables has been proved by their matrix realization. Apply Proposition~\ref{prop:criterion} and Lemma~\ref{lem:positive} to obtain effective associative lifts. The same basis-output positions are nonzero in the two tables, as can be read from Lemma~\ref{lem:table}.

For unimodality in the second case, the factors $D$, $\delta^+$, $W_0^+$, $U_i^+$, and $V_i^+$ are symmetric and unimodal. On $[-4,4]$ the factors $\delta^+\lambda_1^+$ and $\delta^+\lambda_2^+$ have sequences
\[
 (4,11,21,31,34,31,21,11,4),\qquad
 (3,5,10,11,14,11,10,5,3),
\]
so they are symmetric and unimodal as well. A product of symmetric unimodal Laurent polynomials with nonnegative coefficients is again symmetric and unimodal: each such polynomial is a nonnegative linear combination of centered interval sums, and the convolution of two interval sums has a symmetric trapezoidal coefficient sequence. Lemma~\ref{lem:table} now proves unimodality for every dimension coefficient, including $G_i^+$ and $S^+$.

The exact counts are obtained by the same finite multiplication rule. For each triple of basis indices and each weight, add $(d-|c|)/2$ to the cancelling-pair count; add one to the supported-zero count precisely when $c=0$ and $d>0$. Unit coefficients are included. The accompanying independent program \emph{Rank-three matrix lifts} constructs the ten basis vectors, solves every product by exact rational linear algebra, and verifies the stated counts. Its certificate records all $132$ nonzero coefficients, the zero-trace positions, and all five dimension tables. The reported counts are finite exact arithmetic; they make no assertion about a minimum over other dimension tables.
\end{proof}

\subsection{Three further lifts and the exact counts}\label{sec:searched-lifts}
The same formulas produce smaller counts. Let $d$ be a polynomial in $D$ and put
\[
 \delta_1^+=\delta_2^+=d,\qquad G_i^+=a_iDd,\qquad S^+=bDd^2.
\]
The three further choices are
\begin{center}
\begin{tabular}{cclll}
\toprule
Choice & $d$ & $(a_1,a_2,b)$ & $\lambda_1^+$ & $\lambda_2^+$\\
\midrule
A & $D(D-3)$ & $(10,15,30)$ & $4D^2-15D+12$ & $3D^2-10D+4$\\
B & $D(D-3)$ & $(14,7,14)$ & $D$ & $3D^2-10D+4$\\
C & $(D-1)(D-2)$ & $(14,7,14)$ & $D$ & $3D^2-10D+4$\\
\bottomrule
\end{tabular}
\end{center}
Their remaining bottom factors are obtained without matrix inversion:
\[
 W_0^+=\frac{a_1a_2}{b}D,\quad
 U_i^+=\frac{a_i(a_i-1)}bD,\quad
 V_i^+=\frac{a_i}{b}(a_iD^2-\lambda_i^+).
\]
For A these are $5D,3D,7D,2D^2+5D-4,6D^2+5D-2$, in the order $W_0,U_1,U_2,V_1,V_2$. For B and C they are $7D,13D,3D,14D^2-D,2D^2+5D-2$.

\begin{proof}[Proof of Theorem~\ref{thm:lifts-main}]
Every displayed basis is nonsingular, so its table is associative. In A the new first cell factor has Laurent sequence $(4,1,6,1,4)$, which dominates that of $D$ with the required parity. Its bottom factor differs from $V_1=5D-4$ by $2D^2$. All other effectivity checks are those already proved for the first lift.

For B the first cell factor equals the trace factor $D$, while the second is unchanged. The new bottom-factor differences for $W_0,U_1,U_2,V_1,V_2$ are respectively
\[
 6D,\quad12D,\quad2D,\quad14D^2-6D+4,\quad10D-6.
\]
All have nonnegative even Laurent coefficients. The differences $G_1^+-G_1$ and $G_2^+-G_2$, on weights $[-3,3]$, are
\[
 (14,38,30,12,30,38,14),\qquad(6,16,10,0,10,16,6).
\]
The difference $S^+-S$, on $[-5,5]$, is
\[
 (14,52,60,16,6,24,6,16,60,52,14).
\]
These give every remaining inequality and parity check. In C, increasing $d$ by $2$ increases $G_i^+$ and $S^+$ coefficientwise by even nonnegative polynomials; the previous inequalities therefore persist. Proposition~\ref{prop:criterion} proves that all three are effective associative lifts.

For C the factors $D,d,W_0^+,U_i^+,V_i^+$ are symmetric and unimodal. In particular $V_1^+$ and $V_2^+$ have sequences $(14,55,82,55,14)$ and $(2,13,20,13,2)$. The product $d\lambda_1^+=dD$ is unimodal, and $d\lambda_2^+$ has the already displayed sequence $(3,5,10,11,14,11,10,5,3)$. The convolution argument above proves unimodality of all $132$ coefficients.

The following table provides the counts by multiplication type. Each entry is the ordered pair (cancelling multiplicities, supported zero-trace weights). ``Cell output'' excludes the bottom output of a same-cell product. Products in both orders are included. The two earlier choices are denoted I and II, with II the unimodal one.
\begin{center}
\small
\begin{tabular}{lrrrrr}
\toprule
Type & I & II & A & B & C\\
\midrule
$\Sigma^2$ & $(672,2)$ & $(1872,2)$ & $(672,2)$ & $(160,2)$ & $(720,2)$\\
$\Sigma$ times a cell & $(3040,16)$ & $(5040,16)$ & $(3040,16)$ & $(2400,16)$ & $(4080,16)$\\
Same-cell output & $(1060,20)$ & $(1672,24)$ & $(724,20)$ & $(388,4)$ & $(664,8)$\\
Same-cell bottom & $(1580,16)$ & $(1580,16)$ & $(1620,16)$ & $(2940,16)$ & $(2940,16)$\\
Cross-cell bottom & $(1600,0)$ & $(1600,0)$ & $(1600,0)$ & $(2400,0)$ & $(2400,0)$\\
\midrule
Total & $(7952,54)$ & $(11764,58)$ & $(7656,54)$ & $(8288,38)$ & $(10804,42)$\\
\bottomrule
\end{tabular}
\end{center}
Each entry is obtained from Lemma~\ref{lem:table} by expanding $D=q+2+q^{-1}$ and summing $(d_j-|c_j|)/2$, while counting a supported zero once when $c_j=0<d_j$. Unit products contribute $(0,0)$. The displayed multiplication rule gives every input to this finite calculation. For example $S$ has weights $[-4,4]$ whereas $S^+$ has weights $[-5,5]$, accounting for the two supported zero weights in the first row. The independently checked matrix-coordinate computation and a separate word-reduction computation agree on all five tables and these counts. This proves the theorem.
\end{proof}

These constructions were found in a finite parameter search, whose bounds and complete output are recorded in the workbench. The verification here establishes their associativity, effectivity, counts, and stated unimodality. It does not turn a minimum in a search window into a minimum over arbitrary dimension rings or over all inverse-semiring lifts.


\section{Why the bottom ideal forces supported zeros}\label{sec:bottom}
\begin{proof}[Proof of Theorem~\ref{thm:bottom-main}]
For every basis vector $b$, write $b\Sigma=h_b\Sigma$, which is possible by the ideal hypothesis. In particular $h_\gamma=g$. Associativity in $(\gamma\gamma)\Sigma=\gamma(\gamma\Sigma)$ gives
\[
 \sum_b c(\gamma,\gamma;b)h_b=g^2.
\]
Support in $\mathcal S$ is a finite-set semiring: sums take unions and products take Minkowski sums. There is no cancellation of support. The term indexed by $b=\gamma$ therefore gives
\[
 \supp c(\gamma,\gamma;\gamma)+\supp g
 \subseteq\supp(g^2).
\]
The trace of $c(\gamma,\gamma;\gamma)$ has nonzero coefficients at weights $3$ and $-3$, so both weights belong to its support. Write $M=\max\supp g$ and $m=\min\supp g$; the trace of $g$ is nonzero, so its support is nonempty. The inclusion yields $3+M\le2M$ and $-3+m\ge2m$. Hence $M\ge3$ and $m\le-3$. The trace support of $4D(D-1)$ ends at $\pm2$, proving that both extremal occupied weights have trace zero.
\end{proof}

\subsection{An unrestricted lift with no bottom supported zeros}\label{sec:unrestricted}
We now prove that the ideal hypothesis marks a genuine boundary for the supported-zero conclusion. This construction changes the zero pattern of coefficients whose traces vanish.

Write $\Gamma=\{\gamma_w^r:r=1,2,\ w\in\{1,2,12,21\}\}$ and $B=\{1,\Sigma\}\cup\Gamma$. Give these basis elements heights
\[
 h_1=0,\qquad h_\Sigma=2,\qquad h_{\gamma_w^r}=|w|-1.
\]
For $x,y\ne1$, let $b(x,y)$ be the largest exponent in the trace coefficient $c_{xy}^{\Sigma}$. All these coefficients are positive throughout the interval $[-b(x,y),b(x,y)]$. The multiplication rule gives
\[
\begin{array}{c|c}
(x,y)&b(x,y)\\ \hline
(\Sigma,\Sigma)&4\\
(\Sigma,\gamma_w^1)\text{ or }(\gamma_w^1,\Sigma)&2+h_{\gamma_w^1}\\
(\Sigma,\gamma_w^2)\text{ or }(\gamma_w^2,\Sigma)&3+h_{\gamma_w^2}\\
(\gamma_u^1,\gamma_v^1)&h_{\gamma_u^1}+h_{\gamma_v^1}+1-\beta(u,v)\\
(\gamma_u^2,\gamma_v^2)&h_{\gamma_u^2}+h_{\gamma_v^2}+1\\
(\gamma_u^1,\gamma_v^2)\text{ or the reverse}&h_x+h_y+1
\end{array}
\]
Here $\beta(u,v)\in\{0,1\}$ is the alternating-triple exponent in Lemma~\ref{lem:table}; it need not be symmetric in $u,v$. Positivity on each entire interval follows from the displayed $W_0,U_i,V_i,G_i,S$: the nontrivial short factors $5D-4$ and $2D^2-5D+4$ have sequences $(5,6,5)$ and $(2,3,6,3,2)$, and the other factors are products of $D$ and $D-1$ with positive constants.

Consequently, for all nonunit $x,y$,
\begin{equation}\label{eq:bottom-radii}
 b(x,y)\le1+h_x+h_y,\qquad
 \max_{z\in\Gamma}\bigl(b(z,y)-h_z\bigr)=1+h_y
 =\max_{z\in\Gamma}\bigl(b(y,z)-h_z\bigr).
\end{equation}
The maxima are attained by choosing $z$ in the opposite cell from $y$; if $y=\Sigma$, choose $z$ in the second cell.

Define lifted coefficients $\widetilde c_{xy}^v$ as follows. Unit products have their exact Kronecker coefficients. For $x,y\ne1$, the coefficient at $\Sigma$ has exactly the support of $c_{xy}^{\Sigma}$ and the same integer amplitudes. At every other output $v\ne\Sigma$, use the support interval
\[
 [-r(x,y;v),r(x,y;v)],\qquad
 r(x,y;v)=3+h_x+h_y-h_v,
\]
and place the trace coefficient $c_{xy}^v$ on it, using old zeros at all remaining weights. These radii are nonnegative. The interval contains the trace support: the only nonzero nonbottom coefficients of nonunit products have the form
$\delta_iD^\alpha\lambda_i^\beta$, of degree at most $2+\alpha+2\beta$, while the word-length identity gives
\[
 3+h_x+h_y-h_v=2+\alpha+2\beta.
\]
Thus this is a well-defined coefficientwise lift to $\mathcal S$.

\begin{proof}[Proof of Theorem~\ref{thm:no-bottom-main}]
Every associativity equation holds on integer amplitudes because the trace algebra is associative. It remains to check the finite occupied supports. A sum of centered intervals is their union, hence has the largest radius, and their Minkowski product adds radii.

A triple involving $1$ is associative by the exact unit coefficients. Let $x,y,z\ne1$. For a nonbottom output $v\ne\Sigma$, every intermediate element $u\in\Gamma$ contributes the same radius to $(xy)z$:
\[
 r(x,y;u)+r(u,z;v)=6+h_x+h_y+h_z-h_v.
\]
The intermediate bottom element contributes at most that radius, since
\[
 b(x,y)+r(\Sigma,z;v)
 \le(1+h_x+h_y)+(5+h_z-h_v).
\]
An intermediate unit can contribute only when $v=z$, and then its radius $3+h_x+h_y$ is smaller. The corresponding calculation for $x(yz)$ has the same largest radius.

At output $\Sigma$, the intermediate cell elements have largest radius
\[
 \max_{u\in\Gamma}\bigl(r(x,y;u)+b(u,z)\bigr)
   =4+h_x+h_y+h_z
\]
by \eqref{eq:bottom-radii}. The intermediate bottom contribution satisfies
\[
 b(x,y)+b(\Sigma,z)\le4+h_x+h_y+h_z.
\]
If the intermediate unit contributes, then $z=\Sigma$, and its radius $3+h_x+h_y$ is again smaller. Reversing the association gives the same largest radius by the other equality in \eqref{eq:bottom-radii}. This proves every support associativity equation and therefore associativity over $\mathcal S$.

The bottom supports were defined to contain precisely the weights with positive trace coefficient. They contain no occupied zero coefficient. In contrast, the unit coefficient of $\Sigma^2$ has support $[-7,7]$ and identically zero trace. It is not the absent coefficient $\tau$, so the span of $\Sigma$ is not an ideal. The construction proves every asserted property.
\end{proof}

The lift has $829$ nonabsent coefficient triples and $6{,}167$ supported zero-trace weights, all outside the bottom output. These finite counts follow by expanding the stated intervals and the $132$ trace coefficients; an explicit finite-set check verifies all $10{,}000$ coefficientwise support associativity equations. They are not minimum claims. No effective parity lift of this particular support table has been established here. The example refutes the unrestricted supported-zero assertion for inverse-semiring lifts; it does not refute a statement restricted to effective dimension rings with a preserved bottom ideal.


\subsection{Where support must move when the bottom is exact}\label{sec:zero-bottom-tradeoff}
The unrestricted construction can remove supported zeros from the bottom, but it cannot remove all of the support enlargement. The next lower bound applies to every lift with exact bottom supports, whether or not it has an effective refinement.

\begin{proposition}\label{prop:zero-bottom-tradeoff}
Let an associative strict-unit algebra on the stated ten-element basis over $\mathcal S$ have the rank-three trace table. Suppose that each coefficient at the output $\Sigma$ has exactly its nonzero trace support. For every $w\in\{1,2,12,21\}$, put $h_w=|w|-1$. Each of the two ordered products
\[
 \gamma_w^1\Sigma,\qquad \Sigma\gamma_w^1
\]
has a supported zero-trace weight at some cell output of degree at least $3+h_w$, and another at some cell output of degree at most $-3-h_w$. The two cell outputs need not be the same. Consequently there are at least sixteen supported zero-trace weights outside the bottom output. If the algebra has an effective refinement, these weights account for at least sixteen cancelling pairs.
\end{proposition}

\begin{proof}
Write $K_{xy}^v\subset\mathbb Z$ for the finite occupied support of the coefficient of $v$ in $xy$. Addition of supports is union, and multiplication is Minkowski addition. Thus coefficientwise associativity is an equality between unions of such sums, and a weight contributed by one intermediate basis vector cannot be removed by another.

Fix $y=\gamma_w^1$, let $i$ be the first letter of $w$, and put $x=\gamma_i^1$. The matrix identity $P_iP_w=DP_w$ gives the trace coefficient
\[
 c_{xy}^y=\delta_1D=-D(D-1)(D-4).
\]
Its Laurent coefficients at both degrees $3$ and $-3$ are $-1$, so $\{-3,3\}\subseteq K_{xy}^y$. The coefficient of $\Sigma$ in $y\Sigma$ is
\[
 c_{y\Sigma}^{\Sigma}=4D^{1+h_w}(D-1).
\]
Every Laurent coefficient on $[-2-h_w,2+h_w]$ is positive. The exact-bottom hypothesis therefore gives
\[
 K_{y\Sigma}^{\Sigma}=[-2-h_w,2+h_w]\cap\mathbb Z,
 \qquad K_{x\Sigma}^{\Sigma}=[-2,2]\cap\mathbb Z.
\]
The intermediate output $y$ in $(xy)\Sigma$ consequently contributes both weights $5+h_w$ and $-5-h_w$ to the final output $\Sigma$.

Consider the same coefficient in $x(y\Sigma)$. The intermediate output $\Sigma$ contributes only weights in $[-4-h_w,4+h_w]$. The intermediate unit contributes no weight, since the strict-unit coefficient of $\Sigma$ in $x\cdot1=x$ is absent. For every cell basis vector $z\in\Gamma$, the trace table gives
\[
 K_{xz}^{\Sigma}\subseteq[-2,2]\cap\mathbb Z.
\]
Indeed $x$ is a length-one first-cell vector. Against a length-one vector, the bottom radius is at most one; against a length-two vector, it is at most two, as follows directly from the bottom factors $W_0,U_1,V_1$ in Lemma~\ref{lem:table}. Exactness of all bottom supports gives the displayed inclusion for the lifted table as well.

Associativity at the weight $5+h_w$ now requires a cell vector $z$ and weights
\[
 r\in K_{y\Sigma}^z,\qquad s\in K_{xz}^{\Sigma},
 \qquad r+s=5+h_w.
\]
Since $s\le2$, one has $r\ge3+h_w$. Applying the same argument at $-5-h_w$ gives a possibly different cell vector and a weight $r\le-3-h_w$ in its coefficient of $y\Sigma$. Every cell-output trace coefficient of $y\Sigma$ vanishes identically. Both occupied weights are therefore supported zeros.

For the reversed product, let $i$ be the last letter of $w$ and put $x=\gamma_i^1$. The identity $P_wP_i=DP_w$ gives $c_{yx}^y=\delta_1D$. Apply the preceding argument to
\[
 \Sigma(yx)=(\Sigma y)x.
\]
The same bottom-radius bound holds for $K_{zx}^{\Sigma}$, so $\Sigma y$ also has the two asserted supported zero weights.

The four words and the two multiplication orders give eight distinct ordered input pairs. In each pair the positive and negative weights just found are distinct. Counting a supported weight separately for each ordered input pair and output basis vector therefore gives at least sixteen slots. In an effective refinement a supported zero has equal, positive sign multiplicities, so each of those slots contains at least one cancelling pair.
\end{proof}

This proposition is a necessary condition, not an effective construction or an impossibility theorem for exact bottom supports. It leaves open whether the lower bound can be attained and whether any effective lift with exact bottom supports exists.


\subsection{An associative support lift that admits no effective lift}
\label{sec:nonliftable}
An occupied support can satisfy every associativity equation while no choice of finite positive multiplicities on it can do so. The obstruction is already present in \citet[Proposition 12 and Corollary 13, \S6.2--6.3; arXiv version, p.\ 15]{Forsberg}. We now realize it on the fixed rank-three trace table.

Write $T=B\setminus\{1\}$ and $a=\Sigma$. Choose an integer $M$ containing every Laurent trace support in $J=[-M,M]\cap\Z$; the displayed multiplication rule permits $M=5$. Keep the exact unit coefficients. For $a^2$ retain only the output $a$, with occupied support $J$ and its original trace. For every other pair $x,y\in T$, occupy every output $z\in T$ on $J$, with amplitude the original trace coefficient $c_{xy}^z$ and old zeros at the other weights. The output $1$ is absent. This contains every trace support because products of two nonunit basis elements have zero first scalar component in the matrix realization, and $a^2$ has only the output $a$.

\begin{proof}[Proof of Theorem~\ref{thm:nonliftable-main}]
First forget the exponents. On $T$ the output relation is
\[
 a*a=\{a\},\qquad x*y=T\quad\text{if }(x,y)\ne(a,a).
\]
It is associative. If all three inputs are $a$, both associations give $\{a\}$. Otherwise both give $T$: a product involving a different input already has all outputs, and multiplying those outputs again still includes a product other than $a*a$.

Restore exponents. Each nonempty coefficient support is the same interval $J$, so each nonempty two-step contribution has support $J+J$. Thus the equality of the output relations proves equality of the Laurent occupied supports in every nonunit triple. Triples involving $1$ are immediate from the strict unit coefficients. The integer traces are associative by the original matrix algebra, so the resulting coefficients are associative in $\mathcal S$.

Suppose an associative effective lift existed. Evaluate its dimension polynomials at $q=1$, obtaining nonnegative real structure constants $d_{xy}^z$, strictly positive precisely at the occupied outputs. Write
\[
 A=d_{aa}^a>0,\qquad \alpha=d_{ab}^a>0
\]
for any $b\in T\setminus\{a\}$. The coefficient at $a$ in $(aa)b$ is $A\alpha$. The same coefficient in $a(ab)$ is
\[
 A\alpha+\sum_{z\in T\setminus\{a\}}d_{ab}^z d_{az}^a
 >A\alpha.
\]
Indeed every summand in this nonempty sum is positive. This contradicts associativity. The obstruction already rules out arbitrary finite positive real multiplicities, hence also nonnegative integral Laurent dimensions and effective parity classes. The stated lift therefore has no effective lift with the same occupied support.
\end{proof}

Coefficientwise surjectivity of $\pi$ consequently does not imply surjectivity on associative algebra structures, even for this particular rank-three trace algebra. The example in Section~\ref{sec:unrestricted} serves a different purpose: it has no bottom supported zeros. No assertion that those two examples coincide, or that the first example has the present nonliftability obstruction, is needed.


\subsection{The smallest basis rank with an obstruction}\label{sec:minimal-basis-rank}
The two-nonunit output relation in Forsberg's obstruction \citep[Corollary 13]{Forsberg} yields a failure of effective lifting already present for a free coefficient module of rank three. Here ``rank'' refers to the number of named basis vectors, including the unit; it is distinct from the rank-three nonstandard algebra, whose displayed basis has ten elements.

\paragraph{Proof of Theorem~\ref{thm:minimal-basis-rank-main}.}
\begin{proof}
First recall explicitly that $\pi$ is surjective on individual coefficients. At a present weight with integer trace $n>0$, choose $n$ copies of the trivial sign; for $n<0$, choose $-n$ copies of $\eps$; and for a present zero choose $1+\eps$. At an absent weight choose zero. A coefficient has finite support, so this constructs an element of $\mathcal E$ with the prescribed image in $\mathcal S$.

Basis rank one with strict unit has its unique unit multiplication table, which lifts. In basis rank two, write the basis as $1,x$. Its entire multiplication table is determined by
\[
 x^2=\alpha\,1+\beta x.
\]
Choose any coefficient lifts $\widetilde\alpha,\widetilde\beta\in\mathcal E$ and retain the strict unit. This table is associative. By bilinearity it suffices to check basis triples; those containing $1$ are immediate, and the only remaining triple satisfies
\[
 (x^2)x=\widetilde\alpha x+
 \widetilde\beta(\widetilde\alpha\,1+\widetilde\beta x)
 =x(x^2),
\]
using commutativity of the coefficient semiring. Hence every basis-rank-two table lifts.

For rank three, let $e$ denote the old supported zero at Laurent weight zero. On the free $\mathcal S$-semimodule with basis $1,a,b$, retain the strict unit and define
\[
 a^2=e a,\qquad ab=ba=e a+e b,\qquad b^2=e a+e b.
\]
All omitted coefficients are the unsupported zero. Since $e^2=e$ and $e+e=e$, the triples involving two $a$'s and one $b$, or one $a$ and two $b$'s, evaluate to $e a+e b$ in either parenthesization. The triples $a,a,a$ and $b,b,b$ are associative by commutativity, and strict-unit triples are immediate. Bilinearity now proves associativity on the whole free semimodule. Reflection sends every displayed $e$ to the ordinary integer zero, so the trace algebra has square-zero ideal spanned by $a,b$, as stated.

Suppose a coefficient refinement existed and apply its dimension map $\eps\mapsto1$. Write its dimension multiplication as
\[
 a^2=Aa,\qquad ab=Ba+Cb,
\]
where $A,B,C\in\mathbb N[q,q^{-1}]$. The absence of other outputs in $a^2$ is forced: an effective coefficient maps to the unsupported zero only when it is zero. Conversely the present coefficients at the $a$ and $b$ outputs of $ab$ force $B\ne0$ and $C\ne0$. At the $a$ coordinate, associativity gives
\[
 AB=AB+CB.
\]
Evaluating at $q=1$ contradicts $B(1)C(1)>0$. Thus there is no refinement. In fact the prescribed supports force these dimension coefficients to be positive constants, but the contradiction does not require that stronger observation.
\end{proof}


\begin{proposition}[A precise semimodule for the gluing]\label{prop:semimodule}
Let $\mathcal L=\mathcal P_{\mathrm{fin}}(\Z\times B)$, with union and empty set, where $B$ is the ten-element basis. The underlying $\mathcal S$-semimodule of each lifted algebra has, over $U\in\mathcal L$, the group $\Z^U$. Inclusions extend coefficients by zero, a support $S\subset\Z$ acts by $S\cdot U=\{(s+j,b):s\in S,(j,b)\in U\}$, and the coefficient actions are convolution. Its multiplication is the natural bilinear action determined by the lifted structure coefficients.
\end{proposition}
\begin{proof}
An element of the free semimodule is a finitely supported list of Laurent coefficients at the basis vectors, with an occupied exponent allowed to have coefficient zero. Thus it is exactly a finite set $U$ and a vector of integers on $U$. Addition and scalar multiplication give the stated union, transport, and convolution formulas. Associativity of the structure coefficients is precisely the additional identity needed for its bilinear multiplication. This is the module form of the abelian-group diagram in \citet[Theorem 8.5]{CFMM}.
\end{proof}

For the five effective lifts of Theorem~\ref{thm:lifts-main}, the span of $\Sigma$ is a subsemimodule and an ideal; the two four-element cell spans need not be subalgebras, because their products have bottom components. Quotienting by the bottom ideal separates their matrix blocks. This describes their gluing as module data, rather than identifying two zeta functions or two arithmetic schemes without a map.

\subsection{A restriction on changes of basis}\label{sec:rescaling-obstruction}
The exact-bottom question also requires more than a change of the bottom basis vector in one of the five effective examples.

\begin{proposition}\label{prop:rescaling-obstruction}
Work over $\mathbb Q(q)$ in any one of the dimension algebras I, II, A, B, or C. Retain the unit, replace each cell basis vector $\gamma$ by a nonzero scalar multiple $s_\gamma\gamma$, and replace $\Sigma$ by any vector which completes these nine retained or rescaled vectors to a basis. The resulting multiplication table cannot have nonnegative Laurent-polynomial coefficients with every bottom support equal to the corresponding nonzero support of the fixed rank-three trace table. In particular, changing only $\Sigma$ cannot produce an effective lift with exact bottom supports from any of these five examples.
\end{proposition}

\begin{proof}
Put $H=\operatorname{span}_{\mathbb Q(q)}(1,\Gamma)$, and write the new bottom vector as
\[
 \Sigma'=t\Sigma+v,\qquad t\in\mathbb Q(q)^\times,\quad v\in H.
\]
The scalar $t$ is nonzero because the new vectors form a basis. Passing to the one-dimensional vector-space quotient by $H$ shows that the bottom coefficient of any old vector is divided by $t$ in the new basis. Hence the two bottom coefficients in the ordered products of $s_xx$ and $s_yy$ are both multiplied by $s_xs_y/t$. Their ratio is unchanged. This argument does not require $H$ to be an ideal.

Take $x=\gamma_1^1$ and $y=\gamma_{12}^1$. By Lemma~\ref{lem:table}, the old dimension coefficients at $\Sigma$ in $xy$ and $yx$ are, respectively,
\[
 \begin{array}{c|cc}
 \text{dimension algebra}&xy&yx\\ \hline
 \mathrm I,\mathrm{II}&3D^2&2D^2+3D\\
 \mathrm A&3D^2&2D^2+5D-4\\
 \mathrm B,\mathrm C&13D^2&14D^2-D.
 \end{array}
\]
As $q\to+\infty$, one has $D=q+2+q^{-1}\sim q$. The ratio of the two entries therefore tends to $3/2$ in I, II, and A, and to $13/14$ in B and C.

The prescribed trace coefficients in these same ordered products are $D^2$ and $5D-4$. Their nonzero Laurent supports are exactly $[-2,2]\cap\mathbb Z$ and $[-1,1]\cap\mathbb Z$. If the new dimension coefficients $f(q)$ and $g(q)$ were nonnegative Laurent polynomials with exactly those supports, their leading coefficients would be positive and
\[
 \frac{f(q)}{g(q)}\sim c q\qquad(q\to+\infty)
\]
for some $c>0$. This contradicts the invariant finite nonzero limit just computed. The contradiction proves the proposition before any parity condition is imposed, and hence also excludes the asserted effective refinements.
\end{proof}

The conclusion concerns this restricted class of changes of basis. It leaves open changes that mix cell vectors or add other basis components to them, as well as dimension algebras unrelated to the five constructions.

\subsection{Adding bottom components to the cell vectors}\label{sec:bottom-shift-obstruction}
For the two examples with the smallest displayed supported-zero counts, a further family of changes of basis can be excluded even when the new bottom line need not be an ideal.

\begin{proposition}\label{prop:bottom-shift-obstruction}
Start with dimension algebra B or C over $\mathbb Q(q)$. Keep the unit, replace each cell vector by
\[
 \gamma\longmapsto\gamma+r_\gamma\Sigma,
 \qquad r_\gamma\in\mathbb Q(q),
\]
and replace $\Sigma$ by any vector completing these nine vectors to a basis. Suppose the resulting structure coefficients are nonnegative Laurent polynomials and give an effective lift of the fixed rank-three trace table. Among the three bottom coefficients of
\[
 (\gamma_1^1)^2,\qquad
 \gamma_1^1\gamma_{12}^1,\qquad
 \gamma_{12}^1\gamma_1^1
\]
in the new basis, at least one has a supported zero above its trace-support interval and at least one has a supported zero below its trace-support interval. In particular there are at least two bottom supported-zero weights, and this class of changes cannot yield exact bottom supports.
\end{proposition}

\begin{proof}
Evaluate at a sufficiently large or sufficiently small positive real $q$, avoiding the finitely many exceptional points of the rational change of basis. All the proposed new structure constants are nonnegative, so the cone generated by the new basis is pointed and closed under multiplication.

In the old algebra put $e=\Sigma/S$, where $\Sigma^2=S\Sigma$. Then $e$ is a nonzero central idempotent and $S(q)>0$. Products between vectors from different old cells lie in $\mathbb Q(q)e$. Adding multiples of $\Sigma$ to the cell vectors preserves this property. One such cross-cell product has a strictly positive new bottom coefficient, because its prescribed trace coefficient is nonzero. Consequently a nonzero scalar multiple of $e$ belongs to the new positive cone. It follows that either $e$ or $-e$ is positive. The latter alternative is impossible: closure under multiplication would also make $e=(-e)^2$ positive, contradicting pointedness. Thus $e$ is positive, and its new bottom coordinate $t$ is strictly positive.

Let $X,Y$ denote the new cell vectors labelled $\gamma_1^1,\gamma_{12}^1$. There are scalars $u_x,u_y$ with
\[
 eX=u_xe,\qquad eY=u_ye.
\]
Since $e,X,Y$ are positive and $e\ne0$, one has $u_x,u_y\ge0$. The quotient by $\mathbb Q(q)e$ retains the relations
\[
 X^2=cX,\qquad XY=cY,\qquad YX=cX,
 \qquad c=dD>0,
\]
where $d=D(D-3)$ in B and $d=(D-1)(D-2)$ in C. These follow directly from $P_1^2=DP_1$, $P_1P_{12}=DP_{12}$, $P_{12}P_1=DP_1$, and $\lambda_1^+=D$.

Write $P,Q,W$ for the new bottom coefficients of $X^2,XY,YX$. Multiplying the three quotient relations by $e$ determines their missing scalar terms. Because $X$ and $Y$ themselves have zero bottom coordinate in the new basis, this gives
\[
 P=t u_x(u_x-c),\qquad
 Q=t u_y(u_x-c),\qquad
 W=t u_x(u_y-c).
\]
All three are strictly positive: their fixed traces are $D,D^2,5D-4$. Hence $u_x>c$ and $u_y>c$. Moreover
\[
 Q-P=t(u_x-c)(u_y-u_x),\qquad
 W-Q=tc(u_y-u_x).
\]
The two differences therefore have the same sign.

The trace supports of $P,Q,W$ are respectively $[-1,1]$, $[-2,2]$, and $[-1,1]$, intersected with $\mathbb Z$. If none of these coefficients had a supported zero above its trace interval, their maximum exponents would be $1,2,1$. Their positive leading coefficients would imply $Q>P$ and $Q>W$ for all sufficiently large positive $q$, contradicting the same-sign identity. Thus at least one coefficient has a supported zero above its trace interval.

Likewise, if none had a supported zero below its trace interval, their minimum exponents would be $-1,-2,-1$. Letting $q\to0^+$ would again give $Q>P$ and $Q>W$, with the same contradiction. The two resulting weights are distinct because one lies above a positive endpoint and the other below a negative endpoint. This proves the lower bound.
\end{proof}

The proposition permits an arbitrary new bottom vector and does not assume that its span is an ideal. It still concerns a restricted family: mixing cell vectors with one another, adding unit components to them, or changing their individual scalar normalizations may escape this argument.


\section{Integral word algebras at every rank}
This section constructs the integral word order of Theorem~\ref{thm:hecke-word-descent} and an explicit type of denominator certificate. The Hecke normalization is that of \citet[Section 3.1, pp.~20--21]{BSM2013}.

\begin{proof}[Proof of Theorem~\ref{thm:hecke-word-descent}]
The four substitutions $v\mapsto v,-v,v^{-1},-v^{-1}$ are distinct automorphisms of $K_v$ fixing $D$. Conversely, $v$ satisfies
\[
 v^4-(D-2)v^2+1=0.
\]
It follows that $K_v/F_D$ is Galois of degree four, with precisely these automorphisms. Also
\[
 v^{-1}=(D-2)v-v^3.
\]
The quartic relation and this identity show that $1,v,v^2,v^3$ span $R_v$ over $S_D$. Their independence follows from $[K_v:F_D]=4$. Hence $R_v$ is free of rank four over $S_D$, and is flat. Moreover $R_v\otimes_{S_D}F_D=K_v$.

There are two commuting semilinear involutions of the Hecke algebra over $K_v$:
\[
 \alpha(v)=-v,\quad \alpha(T_s)=-T_s;
 \qquad
 \beta(v)=v^{-1},\quad \beta(T_s)=T_s^{-1}.
\]
They preserve the quadratic and braid relations. Indeed the quadratic relation gives
\[
 T_s^{-1}=T_s-v+v^{-1},
\]
and inversion of the braid relations gives the corresponding relations for the inverses. The two involutions commute on scalars and on every $T_s$, where both compositions give $-T_s^{-1}$. Apply each involution diagonally to the two tensor factors. The map $\alpha$ negates both $T_s+v^{-1}$ and $T_s-v$, whereas $\beta$ fixes both. Therefore both involutions fix $P_s$, and hence fix every word.

Choose a $K_v$-basis $w_1,\ldots,w_d$ of $A_{K_v}$ from its spanning set of words. For any word $w$, write
\[
 w=\sum_{i=1}^d a_iw_i,\qquad a_i\in K_v.
\]
Applying $\alpha$ and $\beta$ and using uniqueness of these coordinates shows that each $a_i$ is fixed by the Galois group. Thus $a_i\in F_D$. The same words are consequently an $F_D$-basis of $A_{F_D}$, and scalar extension gives $A_{K_v}$. This proves the field isomorphism and the assertion about multiplication coefficients, because a product of two words is again a word.

The standard Hecke basis makes $M_{R_v}$ free of rank $(r!)^2$ over $R_v$, and therefore free of rank $4(r!)^2$ over $S_D$. The module $A_{S_D}$ is an $S_D$-submodule of this finite free module. Since $S_D$ is Noetherian, it is finite, and its embedding proves torsion-freeness. Word concatenation makes it a unital algebra.

The word basis identifies $A_{F_D}$ with $F_D^d$. Tensor the inclusion $A_{S_D}\subseteq A_{F_D}$ with the flat $S_D$-module $R_v$. This gives an injection
\[
 R_v\otimes_{S_D}A_{S_D}
 \longrightarrow R_v\otimes_{S_D}A_{F_D}
 \cong K_v\otimes_{F_D}A_{F_D}
 \cong A_{K_v}.
\]
Its image is exactly the $R_v$-span of the words, namely $A_{R_v}$. This proves the integral isomorphism. It uses flatness and the field calculation; no unramified integral descent assertion is required.

Finally, fix $a\in A_{S_D}$. Left multiplication by $a$ preserves the finite free $S_D$-module $M_{R_v}$. Its characteristic polynomial is monic over $S_D$, and Cayley--Hamilton, applied to the vector $1$, gives a monic polynomial over $S_D$ annihilating $a$. That polynomial also annihilates $L_a$ on $A_{F_D}$. Every eigenvalue of $L_a$ in an algebraic closure of $F_D$ is therefore integral over $S_D$. The coefficients of $\det(X-L_a)$ are elementary symmetric functions of these eigenvalues, so are integral over $S_D$. They belong to $F_D$, and $S_D=\mathbb Z[D]$ is integrally closed in $F_D$. Hence they belong to $S_D$, as claimed.
\end{proof}

The same calculation gives a finite certificate for a localization on which the specified word algebra is free.

\begin{corollary}\label{cor:hecke-word-localization}
Choose an $F_D$-word basis $w_1,\ldots,w_d$ of $A_{F_D}$. Express its vectors in the ambient $S_D$-basis
\[
 \bigl\{v^a(T_x\otimes T_y):0\le a<4,\quad x,y\in\mathfrak S_r\bigr\}.
\]
Let $f(D)\ne0$ be any $d$-by-$d$ minor of their coordinate matrix. Then
\[
 A_{S_D}[1/f]=\bigoplus_{i=1}^d S_D[1/f]w_i.
\]
In particular all multiplication coefficients in this word basis lie in $S_D[1/f]$.
\end{corollary}

\begin{proof}
The columns are independent over $F_D$, so a nonzero minor exists. Every word has coordinates in $S_D$ in the displayed ambient basis and belongs to the $F_D$-span of the selected words. Restrict its coordinate equations to the rows defining $f$. Cramer's rule puts each coefficient in $f^{-1}S_D\subseteq S_D[1/f]$. Conversely, each $w_i$ already lies in $A_{S_D}$. Localizing gives the asserted equality, and independence follows from the original field basis.
\end{proof}

At $r=4$, the theorem specifies a finite torsion-free integral word algebra and proves that its usual Laurent word algebra is its faithful scalar extension. The corollary supplies a denominator certificate once a word basis and a nonzero minor have been computed. These general arguments alone do not compute the minor, prove integral freeness, or construct a positive distinguished basis. The requirements are separate from the rank-four canonical-basis data: \citet[\S5.2]{Mulmuley} computes only the trivial-block basis element there. In particular this theorem does not identify a least coefficient ring for a positive rank-four basis, and it makes no assertion about descent of individual irreducible characters beyond the regular characteristic polynomials stated above.

\subsection{A specified localization at rank four}\label{sec:rank4-word-minor}

We now give the finite data for one localization. The generic dimension is a published input: \citet[\S11.8, pp.~67--68]{BSM2013} gives $\dim_{K_v}A_{K_v}=114$, with simple-module dimensions
\[
 1,2,3,6,5,5,2,3,1.
\]
Their squared dimensions sum to $114$. The calculation below independently proves that the displayed words are independent; the published result supplies the upper bound and hence their spanning property.

List the permutations of $(1,2,3,4)$ lexicographically as $\pi_0,\ldots,\pi_{23}$. Number the standard tensor coordinates by
\[
 24i+j\quad\longleftrightarrow\quad T_{\pi_i}\otimes T_{\pi_j},
 \qquad 0\le i,j<24.
\]
For a string $a_1\cdots a_k$ in $\{1,2,3\}$, the corresponding word is $P_{s_{a_1}}\cdots P_{s_{a_k}}$; the empty string denotes $1$. Table~\ref{tab:rank4-word-minor} specifies $114$ ordered pairs $(w_i,r_i)$. Let
\[
 C(v)=\bigl([w_j]_{r_i}\bigr)_{1\le i,j\le114},\qquad
 g(v)=\det C(v),
\]
where $[w]_{r}$ is the indicated coordinate of the word. Define the Laurent expression
\begin{equation}\label{eq:rank4-norm-minor}
 f(D)=g(v)g(-v)g(v^{-1})g(-v^{-1}),\qquad
 D=v^2+2+v^{-2}.
\end{equation}
Thus the complete finite definition of $f$ is the Hecke multiplication rule together with the table; no expanded polynomial is required.

\begin{proposition}\label{prop:rank4-word-minor}
The expression in \eqref{eq:rank4-norm-minor} is a nonzero polynomial in $\mathbb Z[D]$. The words of Table~\ref{tab:rank4-word-minor} form an $F_D$-basis of $A_{F_D}$, and
\[
 A_{S_D}[1/f]=\bigoplus_{i=1}^{114}S_D[1/f]w_i.
\]
In particular $S_D[1/f]$ is an explicitly specified coefficient ring for this word basis at rank four. This assertion does not require that $f$ be a least denominator or that the word basis be positive.
\end{proposition}

\begin{proof}
All entries of $C(v)$ lie in $R_v$. Exact arithmetic modulo the prime $p=1000003$ gives
\[
\begin{array}{c|rrrr}
v&2&-2&2^{-1}&-2^{-1}\\\hline
g(v)\bmod p&274637&274637&274637&274637
\end{array}
\]
and consequently
\[
 f(25/4)\equiv988208\pmod p.
\]
The rational number $25/4$ is evaluated using the inverse of $4$ modulo $p$. These nonzero residues prove $g\ne0$ and $f\ne0$ as Laurent polynomials. The $114$ word columns are therefore independent over $K_v$. The published upper bound $114$ proves that they are a $K_v$-basis; Theorem~\ref{thm:hecke-word-descent} then makes them an $F_D$-basis.

For completeness, the modular calculation is a finite certificate with no floating-point comparison. On the standard Hecke basis, right multiplication by $T_s$ sends $T_w$ to $T_{ws}$ when length increases, and to $T_{ws}+(v-v^{-1})T_w$ when length decreases. Expanding the two tensor summands defining $P_s$ computes each word column. Restrict to the listed rows and perform elimination in $\mathbb F_p$; the four displayed values are the resulting determinants. The program \emph{Rank-four word coordinates} records the words, rows and all four values. A separate program, \emph{Rank-four minor replay}, reconstructs the same columns by left multiplication in reversed word order, expands the two tensor summands separately, and uses modular Bareiss elimination. It gives the same four determinants and norm. Both programs and their exact output certificates are included in the accompanying workbench. Trial division up to $\lfloor\sqrt p\rfloor$ also checks primality.

The four factors of $f$ are permuted by $v\mapsto-v$ and $v\mapsto v^{-1}$. Hence $f$ lies in the invariant subring $R_v^{\operatorname{Gal}(K_v/F_D)}$. This subring is exactly $S_D$: sign invariance removes all odd powers of $v$, and inversion invariance writes the result as an integer linear combination of $1$ and $v^{2n}+v^{-2n}$. If $z=v^2+v^{-2}=D-2$, these latter expressions are polynomials in $z$, by the recurrence
\[
 a_0=2,\qquad a_1=z,\qquad a_{n+1}=za_n-a_{n-1}.
\]
Thus $f\in\mathbb Z[D]$, and also $R_v\cap F_D=S_D$.

Let $w$ be any word and write $w=\sum_jx_jw_j$ with $x_j\in F_D$. Apply Cramer's rule to its selected coordinate equations over $R_v$. It gives $g(v)x_j\in R_v$. The quotient $f/g$ is the product of the other three conjugate minors and lies in $R_v$. Therefore $fx_j\in R_v\cap F_D=S_D$. Every word consequently belongs to $f^{-1}\sum_jS_Dw_j$. Localizing gives one inclusion in the asserted equality; the reverse inclusion follows because every $w_j$ already belongs to $A_{S_D}$. Independence was proved above.
\end{proof}

\begin{table}[!ht]
\centering\footnotesize
\setlength{\tabcolsep}{3pt}
\begin{tabular}{lr lr lr lr lr lr }
\toprule
word & row & word & row & word & row & word & row & word & row & word & row\\
\midrule
$\varnothing$ & 0 & \texttt{1} & 6 & \texttt{2} & 2 & \texttt{3} & 1 & \texttt{12} & 8 & \texttt{13} & 7\\
\texttt{21} & 12 & \texttt{23} & 3 & \texttt{32} & 4 & \texttt{121} & 14 & \texttt{123} & 9 & \texttt{132} & 10\\
\texttt{212} & 50 & \texttt{213} & 13 & \texttt{232} & 5 & \texttt{321} & 18 & \texttt{323} & 25 & \texttt{1212} & 54\\
\texttt{1213} & 15 & \texttt{1232} & 11 & \texttt{1321} & 20 & \texttt{1323} & 31 & \texttt{2121} & 56 & \texttt{2123} & 51\\
\texttt{2132} & 16 & \texttt{2321} & 19 & \texttt{2323} & 26 & \texttt{3212} & 52 & \texttt{3213} & 33 & \texttt{3232} & 27\\
\texttt{12121} & 60 & \texttt{12123} & 55 & \texttt{12132} & 17 & \texttt{12321} & 21 & \texttt{12323} & 32 & \texttt{13212} & 61\\
\texttt{13213} & 28 & \texttt{13232} & 34 & \texttt{21213} & 57 & \texttt{21232} & 30 & \texttt{21321} & 22 & \texttt{21323} & 35\\
\texttt{23212} & 36 & \texttt{23213} & 37 & \texttt{23232} & 29 & \texttt{32121} & 58 & \texttt{32123} & 53 & \texttt{32132} & 38\\
\texttt{32321} & 59 & \texttt{121213} & 63 & \texttt{121232} & 39 & \texttt{121321} & 23 & \texttt{121323} & 40 & \texttt{123212} & 41\\
\texttt{123213} & 42 & \texttt{123232} & 43 & \texttt{132121} & 66 & \texttt{132123} & 62 & \texttt{132132} & 44 & \texttt{132321} & 64\\
\texttt{212132} & 45 & \texttt{212321} & 65 & \texttt{212323} & 67 & \texttt{213212} & 75 & \texttt{213213} & 76 & \texttt{213232} & 77\\
\texttt{232121} & 78 & \texttt{232123} & 79 & \texttt{232132} & 80 & \texttt{232321} & 100 & \texttt{321213} & 81 & \texttt{321232} & 82\\
\texttt{321321} & 84 & \texttt{321323} & 101 & \texttt{323212} & 152 & \texttt{1212132} & 46 & \texttt{1212321} & 69 & \texttt{1212323} & 68\\
\texttt{1213212} & 83 & \texttt{1213213} & 85 & \texttt{1213232} & 86 & \texttt{1232121} & 87 & \texttt{1232123} & 88 & \texttt{1232132} & 89\\
\texttt{1232321} & 90 & \texttt{1321213} & 91 & \texttt{1321232} & 92 & \texttt{1321321} & 102 & \texttt{1321323} & 103 & \texttt{1323212} & 104\\
\texttt{2121321} & 105 & \texttt{2121323} & 106 & \texttt{2123212} & 107 & \texttt{2123213} & 150 & \texttt{2123232} & 151 & \texttt{2132121} & 108\\
\texttt{2132123} & 153 & \texttt{2132132} & 109 & \texttt{2132321} & 126 & \texttt{2321232} & 154 & \texttt{2321321} & 155 & \texttt{12121321} & 110\\
\texttt{12121323} & 111 & \texttt{12123212} & 112 & \texttt{12123213} & 113 & \texttt{12123232} & 156 & \texttt{12132121} & 157 & \texttt{12132123} & 114\\
\texttt{12132132} & 127 & \texttt{12132321} & 128 & \texttt{12321232} & 158 & \texttt{12321321} & 159 & \texttt{21213212} & 160 & \texttt{121213212} & 161\\
\bottomrule
\end{tabular}
\caption{Ordered word--coordinate pairs defining the rank-four minor. Read the pairs from left to right across each row, then from top to bottom. The empty word is the unit; coordinate indices begin at zero.}
\label{tab:rank4-word-minor}
\end{table}


This certificate supplies one finite localization for the specified integral word order. It leaves open whether a smaller localization suffices, whether the order is free before localization, and whether a distinguished positive basis exists over an appropriate coefficient ring. None of those questions follows from modular nonvanishing or from the published generic dimension.
\clearpage


\section{What a rank-four or point-count realization still needs}\label{sec:open}
Theorem~\ref{thm:hecke-word-descent} now specifies the integral word order at rank four. The positive-basis problem still needs a precise positive cone and a basis, and a canonical cell order requires its own definition. Proposition~\ref{prop:rank4-word-minor} gives a concrete word basis and a nonzero determinant norm, and therefore one coefficient localization. It does not show that this localization is least or that its word basis is positive. A positive basis over a field and a failure of local freeness for one integral order do not determine a smallest localization. Distinct localizations need not be comparable. Even if a localization $\Z[v,v^{-1},1/f(D)]$ is found, its rational functions need not be traces of finite sign-Tate objects. For example $1/D$ is not a Laurent polynomial, since the only units of $\Z[q,q^{-1}]$ are $\pm q^m$. Thus removing denominators or constructing an arithmetic realization is an additional gate. No minimum coefficient ring or minimum dimension ring at rank four is established here.

\begin{proposition}[An obstruction to a fixed Frobenius interpretation]\label{prop:poles}
Let a sequence $(a_s)_{s\ge1}$ have a finite-dimensional graded trace representation
\[
 a_s=\sum_i(-1)^i\tr(F_i^s\mid V_i)
\]
over a characteristic-zero field embedded in $\C$. Then $\sum_{s\ge1}a_st^s$ is rational and all its nonzero finite poles are simple. Consequently a rational Hilbert series with a pole of order at least two cannot itself be the extension-field point-count series of one variety through such a trace formula.
\end{proposition}
\begin{proof}
Triangularize each $F_i$ over an algebraic closure. Its trace at the $s$th power is the sum of the $s$th powers of its eigenvalues with algebraic multiplicities, even when it has Jordan blocks. Therefore the generating function is a finite integral linear combination of $\alpha t/(1-\alpha t)$. Each term has a simple pole, and taking a sum cannot increase pole orders. A Lefschetz point-count formula is of precisely this form, proving the consequence.
\end{proof}

For each single nonnegative integer $k_n(s)$ one can of course take a disjoint union of $k_n(s)$ rational points over $\F_q$. Allowing the variety to change arbitrarily with $s$ therefore makes the bare existence request automatic. It does not turn pole orders of $\sum k_n(s)t^s$ into weights of one Frobenius operator. A useful geometric construction must add a uniform family, functoriality in $s$, and the cohomological operation whose generating series is being measured. That is the remaining arithmetic gate.

\subsection{A fixed finite trace diagram cannot contain every power}\label{sec:fixed-diagram-obstruction}

The following elementary obstruction applies whenever a proposed realization uses a fixed finite family of vector spaces and endomorphisms. It permits arbitrary insertions, including ones that do not commute with those endomorphisms.

\begin{proposition}\label{prop:fixed-diagram-obstruction}
Fix $q>1$, finitely many finite-dimensional complex vector spaces $V_\ell$, endomorphisms $F_\ell$, and scalars $\mu_\ell$. Fix $c\ne0$ and $\eta\in\mathbb C$. There do not exist endomorphisms $A_{k,\ell}$, chosen separately for each $k\ge1$, such that
\[
 \sum_\ell\mu_\ell\operatorname{Tr}(A_{k,\ell}F_\ell^n)
 =\eta+c(q^n+2+q^{-n})^k
 \qquad(n,k\ge1).
\]
In particular the obstruction is unchanged by adding a fixed weight-zero line to every right side.
\end{proposition}

\begin{proof}
For every fixed $k$, the generating function of the left side is
\[
 \sum_\ell\mu_\ell\operatorname{Tr}
 \bigl(A_{k,\ell}zF_\ell(I-zF_\ell)^{-1}\bigr).
\]
It is rational, with finite poles among the inverses of the nonzero eigenvalues of the fixed matrices $F_\ell$. This is a finite set independent of $k$. No commutation hypothesis is used in the matrix geometric series.

The binomial identity
\[
 (t+2+t^{-1})^k=\sum_{j=-k}^k\binom{2k}{k+j}t^j
\]
shows that the proposed right-side generating function is
\[
 \eta\frac{z}{1-z}
 +c\sum_{j=-k}^k\binom{2k}{k+j}
       \frac{q^jz}{1-q^jz}.
\]
It has an actual pole at $z=q^{-k}$: only its $j=k$ summand is singular there, with nonzero coefficient $c$. The fixed weight-zero term is analytic there. Equality of all coefficients would imply equality of the rational functions, so the fixed finite pole set would have to contain the distinct numbers $q^{-k}$ for every $k\ge1$. This is impossible.
\end{proof}

The proof uses fixed pole locations, and does not assert that inserted traces have only simple poles. Plain traces satisfy $\operatorname{Tr}(F^n)=\sum_\lambda m_\lambda\lambda^n$, because the nilpotent parts of Jordan blocks have zero trace. In contrast,
\[
 F=\begin{pmatrix}\lambda&1\\0&\lambda\end{pmatrix},\qquad
 A=\begin{pmatrix}0&0\\1&0\end{pmatrix},\qquad\lambda\ne0,
\]
give $\operatorname{Tr}(AF^n)=n\lambda^{n-1}$ and the double-pole generating function $z/(1-\lambda z)^2$.

A finite support diagram of finite-dimensional complex vector spaces, equipped with compatible Frobenius operators, gives a finite family of this kind. A M\"obius trace on its support poset uses fixed coefficients $\mu_\ell$. Thus neither such a fixed diagram nor the addition of a rational basepoint can realize all the displayed powers simultaneously. Separate realizations whose spaces depend on the word remain possible. The diagram interpretation uses the inverse additive structure of semimodules in \citet[Proposition~6.1]{FGJM}; Frobenius operators are additional data, not consequences of zero adjunction.


The unrestricted bottom-supported-zero assertion is refuted by Theorem~\ref{thm:no-bottom-main}, while the ideal-preserving obstruction is proved. Theorem~\ref{thm:nonliftable-main} also refutes the assertion that every associative $\mathcal S$-lift comes from an effective dimension table. The effective realizability of the particular zero-bottom example, and the minimum number of supported zero weights in either class, remain open. Those questions require the simultaneous nonnegative parity-constrained associativity equations, not merely a count of cancellations in individual coefficients.

\section*{Note on method}
The zero-adjoined semiring and the inverse-additive-monoid theorem are credited to \citet[Definition 3.6 and Proposition 6.1]{FGJM}, before the project's independent development. That development began in a conversation with Gemini 2.5 Pro in the Gemini app and was formalized by Gemini 2.5 Deep Think in Deep Think mode in November 2025. Versions deposited as HI-AI on Zenodo from 6 November 2025 were later withdrawn. A GPT-5-series Pro model in ChatGPT Pro mode at reasoning effort ultra extracted the algebra in March 2026; such models wrote the April--July corpus, and Codex runs at ultra wrote the August--September workbench. The 21--23 November notes were by Claude Opus 5.5 at max, with Claude Fable 5.2 at max and Claude Opus 5.5 at high consults. They are cited only as unpublished project work. The present mathematics and text were produced by GPT-6 Astra Pro; its numerical effort setting is not exposed. Separate model re-proofs and exact computations are recorded, but no human has checked this text. The newly written material is dedicated to the public domain under CC0 1.0. The workbench is \url{\workbenchurl}; the permanent record address is \url{\recordurl}.

\begin{thebibliography}{99}
\bibitem[Caviglia et al.(2026)]{CFMM} Caviglia, E., Faul, P. F., Manuell, G., \& Mesiti, L. (2026). \emph{Clifford semigroups and the monoidal Grothendieck construction} (Version 2). arXiv:2607.12944. \url{https://arxiv.org/abs/2607.12944}.
\bibitem[Faul et al.(2024)]{FGJM} Faul, P. F., Goswami, A., Joubert, G., \& Manuell, G. (2024). \emph{The case for inverse semirings}. arXiv:2411.12318. \url{https://arxiv.org/abs/2411.12318}.
\bibitem[Mulmuley(2008)]{Mulmuley} Mulmuley, K. D. (2008). \emph{Geometric complexity theory VIII: On canonical bases for the nonstandard quantum groups} (Version 2). arXiv:0709.0751. Conjecture 3.4, \S5.1, and Figure 2. \url{https://arxiv.org/abs/0709.0751}.
\bibitem[Milne(2020)]{MilneANT}
Milne, J. S. (2020). \emph{Algebraic number theory} (Version 3.08).
\url{https://www.jmilne.org/math/CourseNotes/ANT.pdf}.
\bibitem[The Stacks Project Authors(2026)]{StacksRR}
The Stacks Project Authors. (2026). \emph{The Stacks Project}, Tag 0B5B,
Section 53.5: Riemann--Roch.
\url{https://stacks.math.columbia.edu/tag/0B5B}.


\bibitem[Forsberg(2017)]{Forsberg}
Forsberg, L. (2017). Multisemigroups with multiplicities and complete ordered semi-rings. \emph{Beitr\"age zur Algebra und Geometrie, 58}, 405--426. \url{https://doi.org/10.1007/s13366-016-0320-8}. Preprint arXiv:1510.01478, Proposition 12 and Corollary 13, p.\ 15.

\bibitem[Chen and Moriwaki(2019)]{ChenMoriwakiAdelic}
Chen, H., \& Moriwaki, A. (2019). \emph{Arakelov geometry over adelic curves}
(Version 1.2, October 5). Author-hosted manuscript.
\url{https://webusers.imj-prg.fr/~huayi.chen/Recherche/adelic_curve.pdf}.

\bibitem[Serre(1977)]{SerreCharacters}
Serre, J.-P. (1977). \emph{Linear representations of finite groups}
(L. L. Scott, Trans.). Springer.
\url{https://doi.org/10.1007/978-1-4684-9458-7}.
\S2.3, Theorem~4 and Corollary~2, p.~16; \S2.5, Theorem~6, p.~19.

\bibitem[Blasiak et al.(2013)]{BSM2013}
Blasiak, J., Mulmuley, K. D., \& Sohoni, M. (2013).
\emph{Geometric complexity theory IV: Nonstandard quantum group for the Kronecker problem} (Version 4).
arXiv:cs/0703110. Section 3 and equation (22), pp.\ 20--21; Definition 11.1 and equation (93), p.\ 59.
\url{https://arxiv.org/abs/cs/0703110v4}.

\end{thebibliography}
\end{document}
